\documentclass[11pt]{article}

    \usepackage[T2A]{fontenc}
\usepackage[utf8]{inputenc}
\usepackage[english, russian]{babel}

\usepackage[breakable]{tcolorbox}
    \usepackage{parskip} % Stop auto-indenting (to mimic markdown behaviour)
    
    \usepackage{iftex}
    \ifPDFTeX
    	\usepackage[T1]{fontenc}
    	\usepackage{mathpazo}
    \else
    	\usepackage{fontspec}
    \fi

    % Basic figure setup, for now with no caption control since it's done
    % automatically by Pandoc (which extracts ![](path) syntax from Markdown).
    \usepackage{graphicx}
    % Maintain compatibility with old templates. Remove in nbconvert 6.0
    \let\Oldincludegraphics\includegraphics
    % Ensure that by default, figures have no caption (until we provide a
    % proper Figure object with a Caption API and a way to capture that
    % in the conversion process - todo).
    \usepackage{caption}
    \DeclareCaptionFormat{nocaption}{}
    \captionsetup{format=nocaption,aboveskip=0pt,belowskip=0pt}

    \usepackage{float}
    \floatplacement{figure}{H} % forces figures to be placed at the correct location
    \usepackage{xcolor} % Allow colors to be defined
    \usepackage{enumerate} % Needed for markdown enumerations to work
    \usepackage{geometry} % Used to adjust the document margins
    \usepackage{amsmath} % Equations
    \usepackage{amssymb} % Equations
    \usepackage{textcomp} % defines textquotesingle
    % Hack from http://tex.stackexchange.com/a/47451/13684:
    \AtBeginDocument{%
        \def\PYZsq{\textquotesingle}% Upright quotes in Pygmentized code
    }
    \usepackage{upquote} % Upright quotes for verbatim code
    \usepackage{eurosym} % defines \euro
    \usepackage[mathletters]{ucs} % Extended unicode (utf-8) support
    \usepackage{fancyvrb} % verbatim replacement that allows latex
    \usepackage{grffile} % extends the file name processing of package graphics 
                         % to support a larger range
    \makeatletter % fix for old versions of grffile with XeLaTeX
    \@ifpackagelater{grffile}{2019/11/01}
    {
      % Do nothing on new versions
    }
    {
      \def\Gread@@xetex#1{%
        \IfFileExists{"\Gin@base".bb}%
        {\Gread@eps{\Gin@base.bb}}%
        {\Gread@@xetex@aux#1}%
      }
    }
    \makeatother
    \usepackage[Export]{adjustbox} % Used to constrain images to a maximum size
    \adjustboxset{max size={0.9\linewidth}{0.9\paperheight}}

    % The hyperref package gives us a pdf with properly built
    % internal navigation ('pdf bookmarks' for the table of contents,
    % internal cross-reference links, web links for URLs, etc.)
    \usepackage{hyperref}
    % The default LaTeX title has an obnoxious amount of whitespace. By default,
    % titling removes some of it. It also provides customization options.
    \usepackage{titling}
    \usepackage{longtable} % longtable support required by pandoc >1.10
    \usepackage{booktabs}  % table support for pandoc > 1.12.2
    \usepackage[inline]{enumitem} % IRkernel/repr support (it uses the enumerate* environment)
    \usepackage[normalem]{ulem} % ulem is needed to support strikethroughs (\sout)
                                % normalem makes italics be italics, not underlines
    \usepackage{mathrsfs}
    
\usepackage{wrapfig}
\usepackage[rightcaption]{sidecap}
\providecommand{\keywords}[1]{\textbf{\textit{Keywords:}} #1}

\author{A. Yu. Drozdov}

    
    % Colors for the hyperref package
    \definecolor{urlcolor}{rgb}{0,.145,.698}
    \definecolor{linkcolor}{rgb}{.71,0.21,0.01}
    \definecolor{citecolor}{rgb}{.12,.54,.11}

    % ANSI colors
    \definecolor{ansi-black}{HTML}{3E424D}
    \definecolor{ansi-black-intense}{HTML}{282C36}
    \definecolor{ansi-red}{HTML}{E75C58}
    \definecolor{ansi-red-intense}{HTML}{B22B31}
    \definecolor{ansi-green}{HTML}{00A250}
    \definecolor{ansi-green-intense}{HTML}{007427}
    \definecolor{ansi-yellow}{HTML}{DDB62B}
    \definecolor{ansi-yellow-intense}{HTML}{B27D12}
    \definecolor{ansi-blue}{HTML}{208FFB}
    \definecolor{ansi-blue-intense}{HTML}{0065CA}
    \definecolor{ansi-magenta}{HTML}{D160C4}
    \definecolor{ansi-magenta-intense}{HTML}{A03196}
    \definecolor{ansi-cyan}{HTML}{60C6C8}
    \definecolor{ansi-cyan-intense}{HTML}{258F8F}
    \definecolor{ansi-white}{HTML}{C5C1B4}
    \definecolor{ansi-white-intense}{HTML}{A1A6B2}
    \definecolor{ansi-default-inverse-fg}{HTML}{FFFFFF}
    \definecolor{ansi-default-inverse-bg}{HTML}{000000}

    % common color for the border for error outputs.
    \definecolor{outerrorbackground}{HTML}{FFDFDF}

    % commands and environments needed by pandoc snippets
    % extracted from the output of `pandoc -s`
    \providecommand{\tightlist}{%
      \setlength{\itemsep}{0pt}\setlength{\parskip}{0pt}}
    \DefineVerbatimEnvironment{Highlighting}{Verbatim}{commandchars=\\\{\}}
    % Add ',fontsize=\small' for more characters per line
    \newenvironment{Shaded}{}{}
    \newcommand{\KeywordTok}[1]{\textcolor[rgb]{0.00,0.44,0.13}{\textbf{{#1}}}}
    \newcommand{\DataTypeTok}[1]{\textcolor[rgb]{0.56,0.13,0.00}{{#1}}}
    \newcommand{\DecValTok}[1]{\textcolor[rgb]{0.25,0.63,0.44}{{#1}}}
    \newcommand{\BaseNTok}[1]{\textcolor[rgb]{0.25,0.63,0.44}{{#1}}}
    \newcommand{\FloatTok}[1]{\textcolor[rgb]{0.25,0.63,0.44}{{#1}}}
    \newcommand{\CharTok}[1]{\textcolor[rgb]{0.25,0.44,0.63}{{#1}}}
    \newcommand{\StringTok}[1]{\textcolor[rgb]{0.25,0.44,0.63}{{#1}}}
    \newcommand{\CommentTok}[1]{\textcolor[rgb]{0.38,0.63,0.69}{\textit{{#1}}}}
    \newcommand{\OtherTok}[1]{\textcolor[rgb]{0.00,0.44,0.13}{{#1}}}
    \newcommand{\AlertTok}[1]{\textcolor[rgb]{1.00,0.00,0.00}{\textbf{{#1}}}}
    \newcommand{\FunctionTok}[1]{\textcolor[rgb]{0.02,0.16,0.49}{{#1}}}
    \newcommand{\RegionMarkerTok}[1]{{#1}}
    \newcommand{\ErrorTok}[1]{\textcolor[rgb]{1.00,0.00,0.00}{\textbf{{#1}}}}
    \newcommand{\NormalTok}[1]{{#1}}
    
    % Additional commands for more recent versions of Pandoc
    \newcommand{\ConstantTok}[1]{\textcolor[rgb]{0.53,0.00,0.00}{{#1}}}
    \newcommand{\SpecialCharTok}[1]{\textcolor[rgb]{0.25,0.44,0.63}{{#1}}}
    \newcommand{\VerbatimStringTok}[1]{\textcolor[rgb]{0.25,0.44,0.63}{{#1}}}
    \newcommand{\SpecialStringTok}[1]{\textcolor[rgb]{0.73,0.40,0.53}{{#1}}}
    \newcommand{\ImportTok}[1]{{#1}}
    \newcommand{\DocumentationTok}[1]{\textcolor[rgb]{0.73,0.13,0.13}{\textit{{#1}}}}
    \newcommand{\AnnotationTok}[1]{\textcolor[rgb]{0.38,0.63,0.69}{\textbf{\textit{{#1}}}}}
    \newcommand{\CommentVarTok}[1]{\textcolor[rgb]{0.38,0.63,0.69}{\textbf{\textit{{#1}}}}}
    \newcommand{\VariableTok}[1]{\textcolor[rgb]{0.10,0.09,0.49}{{#1}}}
    \newcommand{\ControlFlowTok}[1]{\textcolor[rgb]{0.00,0.44,0.13}{\textbf{{#1}}}}
    \newcommand{\OperatorTok}[1]{\textcolor[rgb]{0.40,0.40,0.40}{{#1}}}
    \newcommand{\BuiltInTok}[1]{{#1}}
    \newcommand{\ExtensionTok}[1]{{#1}}
    \newcommand{\PreprocessorTok}[1]{\textcolor[rgb]{0.74,0.48,0.00}{{#1}}}
    \newcommand{\AttributeTok}[1]{\textcolor[rgb]{0.49,0.56,0.16}{{#1}}}
    \newcommand{\InformationTok}[1]{\textcolor[rgb]{0.38,0.63,0.69}{\textbf{\textit{{#1}}}}}
    \newcommand{\WarningTok}[1]{\textcolor[rgb]{0.38,0.63,0.69}{\textbf{\textit{{#1}}}}}
    
    
    % Define a nice break command that doesn't care if a line doesn't already
    % exist.
    \def\br{\hspace*{\fill} \\* }
    % Math Jax compatibility definitions
    \def\gt{>}
    \def\lt{<}
    \let\Oldtex\TeX
    \let\Oldlatex\LaTeX
    \renewcommand{\TeX}{\textrm{\Oldtex}}
    \renewcommand{\LaTeX}{\textrm{\Oldlatex}}
    % Document parameters
    % Document title
    \title{06.fourth\_Maxwell\_equation}
    
    
    
    
    
% Pygments definitions
\makeatletter
\def\PY@reset{\let\PY@it=\relax \let\PY@bf=\relax%
    \let\PY@ul=\relax \let\PY@tc=\relax%
    \let\PY@bc=\relax \let\PY@ff=\relax}
\def\PY@tok#1{\csname PY@tok@#1\endcsname}
\def\PY@toks#1+{\ifx\relax#1\empty\else%
    \PY@tok{#1}\expandafter\PY@toks\fi}
\def\PY@do#1{\PY@bc{\PY@tc{\PY@ul{%
    \PY@it{\PY@bf{\PY@ff{#1}}}}}}}
\def\PY#1#2{\PY@reset\PY@toks#1+\relax+\PY@do{#2}}

\@namedef{PY@tok@w}{\def\PY@tc##1{\textcolor[rgb]{0.73,0.73,0.73}{##1}}}
\@namedef{PY@tok@c}{\let\PY@it=\textit\def\PY@tc##1{\textcolor[rgb]{0.25,0.50,0.50}{##1}}}
\@namedef{PY@tok@cp}{\def\PY@tc##1{\textcolor[rgb]{0.74,0.48,0.00}{##1}}}
\@namedef{PY@tok@k}{\let\PY@bf=\textbf\def\PY@tc##1{\textcolor[rgb]{0.00,0.50,0.00}{##1}}}
\@namedef{PY@tok@kp}{\def\PY@tc##1{\textcolor[rgb]{0.00,0.50,0.00}{##1}}}
\@namedef{PY@tok@kt}{\def\PY@tc##1{\textcolor[rgb]{0.69,0.00,0.25}{##1}}}
\@namedef{PY@tok@o}{\def\PY@tc##1{\textcolor[rgb]{0.40,0.40,0.40}{##1}}}
\@namedef{PY@tok@ow}{\let\PY@bf=\textbf\def\PY@tc##1{\textcolor[rgb]{0.67,0.13,1.00}{##1}}}
\@namedef{PY@tok@nb}{\def\PY@tc##1{\textcolor[rgb]{0.00,0.50,0.00}{##1}}}
\@namedef{PY@tok@nf}{\def\PY@tc##1{\textcolor[rgb]{0.00,0.00,1.00}{##1}}}
\@namedef{PY@tok@nc}{\let\PY@bf=\textbf\def\PY@tc##1{\textcolor[rgb]{0.00,0.00,1.00}{##1}}}
\@namedef{PY@tok@nn}{\let\PY@bf=\textbf\def\PY@tc##1{\textcolor[rgb]{0.00,0.00,1.00}{##1}}}
\@namedef{PY@tok@ne}{\let\PY@bf=\textbf\def\PY@tc##1{\textcolor[rgb]{0.82,0.25,0.23}{##1}}}
\@namedef{PY@tok@nv}{\def\PY@tc##1{\textcolor[rgb]{0.10,0.09,0.49}{##1}}}
\@namedef{PY@tok@no}{\def\PY@tc##1{\textcolor[rgb]{0.53,0.00,0.00}{##1}}}
\@namedef{PY@tok@nl}{\def\PY@tc##1{\textcolor[rgb]{0.63,0.63,0.00}{##1}}}
\@namedef{PY@tok@ni}{\let\PY@bf=\textbf\def\PY@tc##1{\textcolor[rgb]{0.60,0.60,0.60}{##1}}}
\@namedef{PY@tok@na}{\def\PY@tc##1{\textcolor[rgb]{0.49,0.56,0.16}{##1}}}
\@namedef{PY@tok@nt}{\let\PY@bf=\textbf\def\PY@tc##1{\textcolor[rgb]{0.00,0.50,0.00}{##1}}}
\@namedef{PY@tok@nd}{\def\PY@tc##1{\textcolor[rgb]{0.67,0.13,1.00}{##1}}}
\@namedef{PY@tok@s}{\def\PY@tc##1{\textcolor[rgb]{0.73,0.13,0.13}{##1}}}
\@namedef{PY@tok@sd}{\let\PY@it=\textit\def\PY@tc##1{\textcolor[rgb]{0.73,0.13,0.13}{##1}}}
\@namedef{PY@tok@si}{\let\PY@bf=\textbf\def\PY@tc##1{\textcolor[rgb]{0.73,0.40,0.53}{##1}}}
\@namedef{PY@tok@se}{\let\PY@bf=\textbf\def\PY@tc##1{\textcolor[rgb]{0.73,0.40,0.13}{##1}}}
\@namedef{PY@tok@sr}{\def\PY@tc##1{\textcolor[rgb]{0.73,0.40,0.53}{##1}}}
\@namedef{PY@tok@ss}{\def\PY@tc##1{\textcolor[rgb]{0.10,0.09,0.49}{##1}}}
\@namedef{PY@tok@sx}{\def\PY@tc##1{\textcolor[rgb]{0.00,0.50,0.00}{##1}}}
\@namedef{PY@tok@m}{\def\PY@tc##1{\textcolor[rgb]{0.40,0.40,0.40}{##1}}}
\@namedef{PY@tok@gh}{\let\PY@bf=\textbf\def\PY@tc##1{\textcolor[rgb]{0.00,0.00,0.50}{##1}}}
\@namedef{PY@tok@gu}{\let\PY@bf=\textbf\def\PY@tc##1{\textcolor[rgb]{0.50,0.00,0.50}{##1}}}
\@namedef{PY@tok@gd}{\def\PY@tc##1{\textcolor[rgb]{0.63,0.00,0.00}{##1}}}
\@namedef{PY@tok@gi}{\def\PY@tc##1{\textcolor[rgb]{0.00,0.63,0.00}{##1}}}
\@namedef{PY@tok@gr}{\def\PY@tc##1{\textcolor[rgb]{1.00,0.00,0.00}{##1}}}
\@namedef{PY@tok@ge}{\let\PY@it=\textit}
\@namedef{PY@tok@gs}{\let\PY@bf=\textbf}
\@namedef{PY@tok@gp}{\let\PY@bf=\textbf\def\PY@tc##1{\textcolor[rgb]{0.00,0.00,0.50}{##1}}}
\@namedef{PY@tok@go}{\def\PY@tc##1{\textcolor[rgb]{0.53,0.53,0.53}{##1}}}
\@namedef{PY@tok@gt}{\def\PY@tc##1{\textcolor[rgb]{0.00,0.27,0.87}{##1}}}
\@namedef{PY@tok@err}{\def\PY@bc##1{{\setlength{\fboxsep}{\string -\fboxrule}\fcolorbox[rgb]{1.00,0.00,0.00}{1,1,1}{\strut ##1}}}}
\@namedef{PY@tok@kc}{\let\PY@bf=\textbf\def\PY@tc##1{\textcolor[rgb]{0.00,0.50,0.00}{##1}}}
\@namedef{PY@tok@kd}{\let\PY@bf=\textbf\def\PY@tc##1{\textcolor[rgb]{0.00,0.50,0.00}{##1}}}
\@namedef{PY@tok@kn}{\let\PY@bf=\textbf\def\PY@tc##1{\textcolor[rgb]{0.00,0.50,0.00}{##1}}}
\@namedef{PY@tok@kr}{\let\PY@bf=\textbf\def\PY@tc##1{\textcolor[rgb]{0.00,0.50,0.00}{##1}}}
\@namedef{PY@tok@bp}{\def\PY@tc##1{\textcolor[rgb]{0.00,0.50,0.00}{##1}}}
\@namedef{PY@tok@fm}{\def\PY@tc##1{\textcolor[rgb]{0.00,0.00,1.00}{##1}}}
\@namedef{PY@tok@vc}{\def\PY@tc##1{\textcolor[rgb]{0.10,0.09,0.49}{##1}}}
\@namedef{PY@tok@vg}{\def\PY@tc##1{\textcolor[rgb]{0.10,0.09,0.49}{##1}}}
\@namedef{PY@tok@vi}{\def\PY@tc##1{\textcolor[rgb]{0.10,0.09,0.49}{##1}}}
\@namedef{PY@tok@vm}{\def\PY@tc##1{\textcolor[rgb]{0.10,0.09,0.49}{##1}}}
\@namedef{PY@tok@sa}{\def\PY@tc##1{\textcolor[rgb]{0.73,0.13,0.13}{##1}}}
\@namedef{PY@tok@sb}{\def\PY@tc##1{\textcolor[rgb]{0.73,0.13,0.13}{##1}}}
\@namedef{PY@tok@sc}{\def\PY@tc##1{\textcolor[rgb]{0.73,0.13,0.13}{##1}}}
\@namedef{PY@tok@dl}{\def\PY@tc##1{\textcolor[rgb]{0.73,0.13,0.13}{##1}}}
\@namedef{PY@tok@s2}{\def\PY@tc##1{\textcolor[rgb]{0.73,0.13,0.13}{##1}}}
\@namedef{PY@tok@sh}{\def\PY@tc##1{\textcolor[rgb]{0.73,0.13,0.13}{##1}}}
\@namedef{PY@tok@s1}{\def\PY@tc##1{\textcolor[rgb]{0.73,0.13,0.13}{##1}}}
\@namedef{PY@tok@mb}{\def\PY@tc##1{\textcolor[rgb]{0.40,0.40,0.40}{##1}}}
\@namedef{PY@tok@mf}{\def\PY@tc##1{\textcolor[rgb]{0.40,0.40,0.40}{##1}}}
\@namedef{PY@tok@mh}{\def\PY@tc##1{\textcolor[rgb]{0.40,0.40,0.40}{##1}}}
\@namedef{PY@tok@mi}{\def\PY@tc##1{\textcolor[rgb]{0.40,0.40,0.40}{##1}}}
\@namedef{PY@tok@il}{\def\PY@tc##1{\textcolor[rgb]{0.40,0.40,0.40}{##1}}}
\@namedef{PY@tok@mo}{\def\PY@tc##1{\textcolor[rgb]{0.40,0.40,0.40}{##1}}}
\@namedef{PY@tok@ch}{\let\PY@it=\textit\def\PY@tc##1{\textcolor[rgb]{0.25,0.50,0.50}{##1}}}
\@namedef{PY@tok@cm}{\let\PY@it=\textit\def\PY@tc##1{\textcolor[rgb]{0.25,0.50,0.50}{##1}}}
\@namedef{PY@tok@cpf}{\let\PY@it=\textit\def\PY@tc##1{\textcolor[rgb]{0.25,0.50,0.50}{##1}}}
\@namedef{PY@tok@c1}{\let\PY@it=\textit\def\PY@tc##1{\textcolor[rgb]{0.25,0.50,0.50}{##1}}}
\@namedef{PY@tok@cs}{\let\PY@it=\textit\def\PY@tc##1{\textcolor[rgb]{0.25,0.50,0.50}{##1}}}

\def\PYZbs{\char`\\}
\def\PYZus{\char`\_}
\def\PYZob{\char`\{}
\def\PYZcb{\char`\}}
\def\PYZca{\char`\^}
\def\PYZam{\char`\&}
\def\PYZlt{\char`\<}
\def\PYZgt{\char`\>}
\def\PYZsh{\char`\#}
\def\PYZpc{\char`\%}
\def\PYZdl{\char`\$}
\def\PYZhy{\char`\-}
\def\PYZsq{\char`\'}
\def\PYZdq{\char`\"}
\def\PYZti{\char`\~}
% for compatibility with earlier versions
\def\PYZat{@}
\def\PYZlb{[}
\def\PYZrb{]}
\makeatother


    % For linebreaks inside Verbatim environment from package fancyvrb. 
    \makeatletter
        \newbox\Wrappedcontinuationbox 
        \newbox\Wrappedvisiblespacebox 
        \newcommand*\Wrappedvisiblespace {\textcolor{red}{\textvisiblespace}} 
        \newcommand*\Wrappedcontinuationsymbol {\textcolor{red}{\llap{\tiny$\m@th\hookrightarrow$}}} 
        \newcommand*\Wrappedcontinuationindent {3ex } 
        \newcommand*\Wrappedafterbreak {\kern\Wrappedcontinuationindent\copy\Wrappedcontinuationbox} 
        % Take advantage of the already applied Pygments mark-up to insert 
        % potential linebreaks for TeX processing. 
        %        {, <, #, %, $, ' and ": go to next line. 
        %        _, }, ^, &, >, - and ~: stay at end of broken line. 
        % Use of \textquotesingle for straight quote. 
        \newcommand*\Wrappedbreaksatspecials {% 
            \def\PYGZus{\discretionary{\char`\_}{\Wrappedafterbreak}{\char`\_}}% 
            \def\PYGZob{\discretionary{}{\Wrappedafterbreak\char`\{}{\char`\{}}% 
            \def\PYGZcb{\discretionary{\char`\}}{\Wrappedafterbreak}{\char`\}}}% 
            \def\PYGZca{\discretionary{\char`\^}{\Wrappedafterbreak}{\char`\^}}% 
            \def\PYGZam{\discretionary{\char`\&}{\Wrappedafterbreak}{\char`\&}}% 
            \def\PYGZlt{\discretionary{}{\Wrappedafterbreak\char`\<}{\char`\<}}% 
            \def\PYGZgt{\discretionary{\char`\>}{\Wrappedafterbreak}{\char`\>}}% 
            \def\PYGZsh{\discretionary{}{\Wrappedafterbreak\char`\#}{\char`\#}}% 
            \def\PYGZpc{\discretionary{}{\Wrappedafterbreak\char`\%}{\char`\%}}% 
            \def\PYGZdl{\discretionary{}{\Wrappedafterbreak\char`\$}{\char`\$}}% 
            \def\PYGZhy{\discretionary{\char`\-}{\Wrappedafterbreak}{\char`\-}}% 
            \def\PYGZsq{\discretionary{}{\Wrappedafterbreak\textquotesingle}{\textquotesingle}}% 
            \def\PYGZdq{\discretionary{}{\Wrappedafterbreak\char`\"}{\char`\"}}% 
            \def\PYGZti{\discretionary{\char`\~}{\Wrappedafterbreak}{\char`\~}}% 
        } 
        % Some characters . , ; ? ! / are not pygmentized. 
        % This macro makes them "active" and they will insert potential linebreaks 
        \newcommand*\Wrappedbreaksatpunct {% 
            \lccode`\~`\.\lowercase{\def~}{\discretionary{\hbox{\char`\.}}{\Wrappedafterbreak}{\hbox{\char`\.}}}% 
            \lccode`\~`\,\lowercase{\def~}{\discretionary{\hbox{\char`\,}}{\Wrappedafterbreak}{\hbox{\char`\,}}}% 
            \lccode`\~`\;\lowercase{\def~}{\discretionary{\hbox{\char`\;}}{\Wrappedafterbreak}{\hbox{\char`\;}}}% 
            \lccode`\~`\:\lowercase{\def~}{\discretionary{\hbox{\char`\:}}{\Wrappedafterbreak}{\hbox{\char`\:}}}% 
            \lccode`\~`\?\lowercase{\def~}{\discretionary{\hbox{\char`\?}}{\Wrappedafterbreak}{\hbox{\char`\?}}}% 
            \lccode`\~`\!\lowercase{\def~}{\discretionary{\hbox{\char`\!}}{\Wrappedafterbreak}{\hbox{\char`\!}}}% 
            \lccode`\~`\/\lowercase{\def~}{\discretionary{\hbox{\char`\/}}{\Wrappedafterbreak}{\hbox{\char`\/}}}% 
            \catcode`\.\active
            \catcode`\,\active 
            \catcode`\;\active
            \catcode`\:\active
            \catcode`\?\active
            \catcode`\!\active
            \catcode`\/\active 
            \lccode`\~`\~ 	
        }
    \makeatother

    \let\OriginalVerbatim=\Verbatim
    \makeatletter
    \renewcommand{\Verbatim}[1][1]{%
        %\parskip\z@skip
        \sbox\Wrappedcontinuationbox {\Wrappedcontinuationsymbol}%
        \sbox\Wrappedvisiblespacebox {\FV@SetupFont\Wrappedvisiblespace}%
        \def\FancyVerbFormatLine ##1{\hsize\linewidth
            \vtop{\raggedright\hyphenpenalty\z@\exhyphenpenalty\z@
                \doublehyphendemerits\z@\finalhyphendemerits\z@
                \strut ##1\strut}%
        }%
        % If the linebreak is at a space, the latter will be displayed as visible
        % space at end of first line, and a continuation symbol starts next line.
        % Stretch/shrink are however usually zero for typewriter font.
        \def\FV@Space {%
            \nobreak\hskip\z@ plus\fontdimen3\font minus\fontdimen4\font
            \discretionary{\copy\Wrappedvisiblespacebox}{\Wrappedafterbreak}
            {\kern\fontdimen2\font}%
        }%
        
        % Allow breaks at special characters using \PYG... macros.
        \Wrappedbreaksatspecials
        % Breaks at punctuation characters . , ; ? ! and / need catcode=\active 	
        \OriginalVerbatim[#1,codes*=\Wrappedbreaksatpunct]%
    }
    \makeatother

    % Exact colors from NB
    \definecolor{incolor}{HTML}{303F9F}
    \definecolor{outcolor}{HTML}{D84315}
    \definecolor{cellborder}{HTML}{CFCFCF}
    \definecolor{cellbackground}{HTML}{F7F7F7}
    
    % prompt
    \makeatletter
    \newcommand{\boxspacing}{\kern\kvtcb@left@rule\kern\kvtcb@boxsep}
    \makeatother
    \newcommand{\prompt}[4]{
        {\ttfamily\llap{{\color{#2}[#3]:\hspace{3pt}#4}}\vspace{-\baselineskip}}
    }
    

    
    % Prevent overflowing lines due to hard-to-break entities
    \sloppy 
    % Setup hyperref package
    \hypersetup{
      breaklinks=true,  % so long urls are correctly broken across lines
      colorlinks=true,
      urlcolor=urlcolor,
      linkcolor=linkcolor,
      citecolor=citecolor,
      }
    % Slightly bigger margins than the latex defaults
    
    \geometry{verbose,tmargin=1in,bmargin=1in,lmargin=1in,rmargin=1in}
    
    

\begin{document}
    
    \maketitle
    
    

    
    \section{Четвертое уравнение
Максвелла}\label{ux447ux435ux442ux432ux435ux440ux442ux43eux435-ux443ux440ux430ux432ux43dux435ux43dux438ux435-ux43cux430ux43aux441ux432ux435ux43bux43bux430}

На форуме sciteclibrary.ru 04.06.18 появилась тема «Как запаздывающий
Лиенар-Вихерт становится "незапаздывающим". Визуализация»
http://www.sciteclibrary.ru/cgi-bin/yabb2/YaBB.pl?num=1528093569/0 Автор
темы, Дробышев, в частности, пишет «Забавно, но в процессе перехода
возникают локальные минимумы скалярного потенциала, например, здесь:

Изображение взято для момента \(t=7{,}5\), когда заряд уже движется
равномерно. Естественно, возникновение потенциальных ям не связано с
какими-то дополнительными зарядами, материализующимися "из воздуха". Это
ведь не статика, где \(\vec{E}=-\nabla \varphi\), а динамика с
\(\vec{E}=-\nabla \varphi -(1/c)\partial\vec{A}/ \partial t\), поэтому с
уравнением \(\nabla\cdot\vec{E}=0\) по-прежнему все в порядке.». Конец
цитаты. Приведенное наблюдение наталкивает на мысль проверить
непосредственным дифференцированием, равна ли нулю дивергенция
электрического поля в вакууме при ускоренном движении одиночного заряда.
Градиент скалярного потенциала ЛВ (первое слагаемое выражения
электрического поля через потенциалы)
\[\overrightarrow{{{E}_{1}}}=-\nabla \varphi =-\nabla \frac{e}{K}=\frac{e}{{{K}^{2}}}\left\{ \frac{\overrightarrow{R}}{K}\left( 1+\frac{\overrightarrow{a}\cdot \overrightarrow{R}}{{{c}^{2}}}-\frac{\overrightarrow{v}\cdot \overrightarrow{v}}{{{c}^{2}}} \right)-\frac{\overrightarrow{v}}{c} \right\}\]
\[\overrightarrow{{{E}_{1}}}=e\frac{\overrightarrow{R}}{{{K}^{3}}}\left( 1+\frac{\overrightarrow{a}\cdot \overrightarrow{R}}{{{c}^{2}}}-\frac{\overrightarrow{v}\cdot \overrightarrow{v}}{{{c}^{2}}} \right)-\frac{e}{{{K}^{2}}}\frac{\overrightarrow{v}}{c}\]
Для расчёта дивергенции выразим производную первого слагаемого
электрического поля по координате точки наблюдения \[\begin{align}
  & \frac{1}{e}\frac{\partial \overrightarrow{{{E}_{1}}}}{\partial x}=-\frac{\overrightarrow{R}}{{{K}^{4}}}\frac{{{R}_{x}}}{{{c}^{2}}}\left( \frac{\overrightarrow{R}}{c}\frac{\partial \overrightarrow{a}}{\partial t'} \right)+\frac{\overrightarrow{R}}{{{K}^{4}}}\frac{{{R}_{x}}}{{{c}^{2}}}\left( \frac{3\ \overrightarrow{v}\cdot \overrightarrow{a}}{c} \right)+\frac{\overrightarrow{R}}{{{K}^{3}}}\frac{{{a}_{x}}}{{{c}^{2}}}+\frac{{{R}_{x}}}{{{K}^{3}}}\frac{\overrightarrow{a}}{{{c}^{2}}}-\frac{2}{{{K}^{3}}}\frac{\overrightarrow{v}}{c}\frac{{{v}_{x}}}{c}+ \\ 
 & +\left( 1+\frac{\overrightarrow{a}\cdot \overrightarrow{R}}{{{c}^{2}}}-\frac{\overrightarrow{v}\cdot \overrightarrow{v}}{{{c}^{2}}} \right)\cdot \left\{ \frac{\overrightarrow{i}}{{{K}^{3}}}-\frac{3\overrightarrow{R}}{K}\frac{{{R}_{x}}}{{{K}^{4}}}\left( 1+\frac{\overrightarrow{R}\cdot \overrightarrow{a}-{{v}^{2}}}{{{c}^{2}}} \right)+\frac{3\overrightarrow{R}}{{{K}^{4}}}\frac{{{v}_{x}}}{c}+\frac{3{{R}_{x}}}{{{K}^{4}}}\frac{\overrightarrow{v}}{c} \right\} \\ 
\end{align}\] Выражения производной первого слагаемого электрического
поля по остальным координатам \(y\) и \(z\) точки наблюдения легко
получить меняя в данном выражении \({{R}_{x}}\) \({{a}_{x}}\)
\({{v}_{x}}\) \(\overrightarrow{i}\) на \({{R}_{y}}\) \({{a}_{y}}\)
\({{v}_{y}}\) \(\overrightarrow{j}\)и на \({{R}_{z}}\) \({{a}_{z}}\)
\({{v}_{z}}\) \(\overrightarrow{k}\)соответственно Возьмём теперь
производную \(x\) компоненты первого слагаемого электрического поля по
\(x\) координате точки наблюдения \[\begin{align}
  & \frac{1}{e}\frac{\partial {{E}_{{{1}_{x}}}}}{\partial x}=-\frac{{{R}_{x}}}{{{K}^{4}}}\frac{{{R}_{x}}}{{{c}^{2}}}\left( \frac{\overrightarrow{R}}{c}\frac{\partial \overrightarrow{a}}{\partial t'} \right)+\frac{{{R}_{x}}}{{{K}^{4}}}\frac{{{R}_{x}}}{{{c}^{2}}}\left( \frac{3\ \overrightarrow{v}\cdot \overrightarrow{a}}{c} \right)+\frac{{{R}_{x}}}{{{K}^{3}}}\frac{{{a}_{x}}}{{{c}^{2}}}+\frac{{{R}_{x}}}{{{K}^{3}}}\frac{{{a}_{x}}}{{{c}^{2}}}-\frac{2}{{{K}^{3}}}\frac{{{v}_{x}}}{c}\frac{{{v}_{x}}}{c}+ \\ 
 & +\left( 1+\frac{\overrightarrow{a}\cdot \overrightarrow{R}}{{{c}^{2}}}-\frac{\overrightarrow{v}\cdot \overrightarrow{v}}{{{c}^{2}}} \right)\cdot \left\{ \frac{1}{{{K}^{3}}}-\frac{3{{R}_{x}}}{K}\frac{{{R}_{x}}}{{{K}^{4}}}\left( 1+\frac{\overrightarrow{R}\cdot \overrightarrow{a}-{{v}^{2}}}{{{c}^{2}}} \right)+\frac{3{{R}_{x}}}{{{K}^{4}}}\frac{{{v}_{x}}}{c}+\frac{3{{R}_{x}}}{{{K}^{4}}}\frac{{{v}_{x}}}{c} \right\} \\ 
\end{align}\] Выражения производных \(y\) и \(z\) компонент первого
слагаемого электрического поля по \(y\)и \(z\) координатам точки
наблюдения легко получить аналогичной заменой индексов в данном
выражении Скалярно складывая производные\(x\) , \(y\) и \(z\)компонент
первого слагаемого электрического поля по соответственно \(x\) , \(y\) и
\(z\)координатам точки наблюдения получаем выражение для дивергенции
первого слагаемого электрического поля \[\begin{align}
  & \frac{1}{e}{{\left( \frac{\partial {{E}_{x}}}{\partial x}+\frac{\partial {{E}_{y}}}{\partial y}+\frac{\partial {{E}_{z}}}{\partial z} \right)}_{1}}=-\frac{{{R}^{2}}}{{{K}^{4}}{{c}^{2}}}\left( \frac{\overrightarrow{R}}{c}\frac{\partial \overrightarrow{a}}{\partial t'} \right)+\frac{{{R}^{2}}}{{{K}^{4}}{{c}^{2}}}\left( \frac{3\ \overrightarrow{v}\cdot \overrightarrow{a}}{c} \right)+\frac{2}{{{K}^{3}}}\frac{\overrightarrow{a}\cdot \overrightarrow{R}}{{{c}^{2}}}-\frac{2}{{{K}^{3}}}\frac{{{v}^{2}}}{{{c}^{2}}}+ \\ 
 & +\left( 1+\frac{\overrightarrow{a}\cdot \overrightarrow{R}}{{{c}^{2}}}-\frac{\overrightarrow{v}\cdot \overrightarrow{v}}{{{c}^{2}}} \right)\cdot \left\{ \frac{3}{{{K}^{3}}}-\frac{3{{R}^{2}}}{{{K}^{5}}}\left( 1+\frac{\overrightarrow{R}\cdot \overrightarrow{a}-{{v}^{2}}}{{{c}^{2}}} \right)+6\frac{\overrightarrow{v}\cdot \overrightarrow{R}}{c{{K}^{4}}} \right\} \\ 
\end{align}\]

\[\begin{align}
  & \frac{1}{e}{{\left( \frac{\partial {{E}_{x}}}{\partial x}+\frac{\partial {{E}_{y}}}{\partial y}+\frac{\partial {{E}_{z}}}{\partial z} \right)}_{1}}=-\frac{{{R}^{2}}}{{{K}^{4}}{{c}^{2}}}\left( \frac{\overrightarrow{R}}{c}\frac{\partial \overrightarrow{a}}{\partial t'} \right)+\frac{{{R}^{2}}}{{{K}^{4}}{{c}^{2}}}\left( \frac{3\ \overrightarrow{v}\cdot \overrightarrow{a}}{c} \right)+\frac{2}{{{K}^{3}}}\left( \frac{\overrightarrow{a}\cdot \overrightarrow{R}}{{{c}^{2}}}-\frac{{{v}^{2}}}{{{c}^{2}}} \right)+ \\ 
 & +\left( 1+\frac{\overrightarrow{a}\cdot \overrightarrow{R}}{{{c}^{2}}}-\frac{\overrightarrow{v}\cdot \overrightarrow{v}}{{{c}^{2}}} \right)\cdot \left\{ \frac{3}{{{K}^{3}}}-\frac{3{{R}^{2}}}{{{K}^{5}}}\left( 1+\frac{\overrightarrow{R}\cdot \overrightarrow{a}-{{v}^{2}}}{{{c}^{2}}} \right)+6\frac{\overrightarrow{v}\cdot \overrightarrow{R}}{c{{K}^{4}}} \right\} \\ 
\end{align}\]

\[\begin{align}
  & \frac{1}{e}{{\left( \frac{\partial {{E}_{x}}}{\partial x}+\frac{\partial {{E}_{y}}}{\partial y}+\frac{\partial {{E}_{z}}}{\partial z} \right)}_{1}}=-\frac{{{R}^{2}}}{{{K}^{4}}{{c}^{2}}}\left( \frac{\overrightarrow{R}}{c}\frac{\partial \overrightarrow{a}}{\partial t'} \right)+\frac{{{R}^{2}}}{{{K}^{4}}{{c}^{2}}}\left( \frac{3\ \overrightarrow{v}\cdot \overrightarrow{a}}{c} \right)+\frac{2}{{{K}^{3}}}\left( 1+\frac{\overrightarrow{a}\cdot \overrightarrow{R}}{{{c}^{2}}}-\frac{{{v}^{2}}}{{{c}^{2}}} \right)-\frac{2}{{{K}^{3}}}+ \\ 
 & +\left( 1+\frac{\overrightarrow{a}\cdot \overrightarrow{R}}{{{c}^{2}}}-\frac{\overrightarrow{v}\cdot \overrightarrow{v}}{{{c}^{2}}} \right)\cdot \left\{ \frac{3}{{{K}^{3}}}-\frac{3{{R}^{2}}}{{{K}^{5}}}\left( 1+\frac{\overrightarrow{R}\cdot \overrightarrow{a}-{{v}^{2}}}{{{c}^{2}}} \right)+6\frac{\overrightarrow{v}\cdot \overrightarrow{R}}{c{{K}^{4}}} \right\} \\ 
\end{align}\]

\[\begin{align}
  & \frac{1}{e}{{\left( \frac{\partial {{E}_{x}}}{\partial x}+\frac{\partial {{E}_{y}}}{\partial y}+\frac{\partial {{E}_{z}}}{\partial z} \right)}_{1}}=-\frac{{{R}^{2}}}{{{K}^{4}}{{c}^{2}}}\left( \frac{\overrightarrow{R}}{c}\frac{\partial \overrightarrow{a}}{\partial t'} \right)+\frac{{{R}^{2}}}{{{K}^{4}}{{c}^{2}}}\left( \frac{3\ \overrightarrow{v}\cdot \overrightarrow{a}}{c} \right)-\frac{2}{{{K}^{3}}}+ \\ 
 & +\left( 1+\frac{\overrightarrow{a}\cdot \overrightarrow{R}}{{{c}^{2}}}-\frac{\overrightarrow{v}\cdot \overrightarrow{v}}{{{c}^{2}}} \right)\cdot \left\{ \frac{5}{{{K}^{3}}}-\frac{3{{R}^{2}}}{{{K}^{5}}}\left( 1+\frac{\overrightarrow{R}\cdot \overrightarrow{a}-{{v}^{2}}}{{{c}^{2}}} \right)+6\frac{\overrightarrow{v}\cdot \overrightarrow{R}}{c{{K}^{4}}} \right\} \\ 
\end{align}\]

Второе слагаемое выражения электрического поля через потенциалы
http://www.sciteclibrary.ru/cgi-bin/yabb2/YaBB.pl?num=1528093569/417\#417
\[\overrightarrow{{{E}_{2}}}=-\frac{1}{c}\frac{\partial \overrightarrow{A}}{\partial t}=\frac{e}{{{K}^{2}}}\left[ \frac{\overrightarrow{v}}{c}\left\{ \frac{R}{K}\left( \frac{{{v}^{2}}}{{{c}^{2}}}-\frac{\overrightarrow{a}\cdot \overrightarrow{R}}{{{c}^{2}}}-1 \right)+1 \right\}-R\frac{\overrightarrow{a}}{{{c}^{2}}} \right]\]
\[\overrightarrow{{{E}_{2}}}=\frac{e}{{{K}^{2}}}\frac{\overrightarrow{v}}{c}\left\{ \frac{R}{K}\left( \frac{{{v}^{2}}}{{{c}^{2}}}-\frac{\overrightarrow{a}\cdot \overrightarrow{R}}{{{c}^{2}}}-1 \right)+1 \right\}-\frac{e}{{{K}^{2}}}R\frac{\overrightarrow{a}}{{{c}^{2}}}\]
\[\overrightarrow{{{E}_{2}}}=\frac{eR}{{{K}^{3}}}\frac{\overrightarrow{v}}{c}\left( \frac{{{v}^{2}}}{{{c}^{2}}}-\frac{\overrightarrow{a}\cdot \overrightarrow{R}}{{{c}^{2}}}-1 \right)+\frac{e}{{{K}^{2}}}\frac{\overrightarrow{v}}{c}-\frac{eR}{{{K}^{2}}}\frac{\overrightarrow{a}}{{{c}^{2}}}\]
\[\frac{1}{e}\frac{\partial \overrightarrow{{{E}_{2}}}}{\partial x}=\frac{\partial }{\partial x}\left[ \frac{R}{{{K}^{3}}}\frac{\overrightarrow{v}}{c}\left( \frac{{{v}^{2}}}{{{c}^{2}}}-\frac{\overrightarrow{a}\cdot \overrightarrow{R}}{{{c}^{2}}}-1 \right)+\frac{1}{{{K}^{2}}}\frac{\overrightarrow{v}}{c}-\frac{R}{{{K}^{2}}}\frac{\overrightarrow{a}}{{{c}^{2}}} \right]\]
\[\frac{1}{e}\frac{\partial \overrightarrow{{{E}_{2}}}}{\partial x}=\left( \frac{R}{{{K}^{3}}}\frac{\overrightarrow{v}}{c} \right)\frac{\partial }{\partial x}\left( \frac{{{v}^{2}}}{{{c}^{2}}}-\frac{\overrightarrow{a}\cdot \overrightarrow{R}}{{{c}^{2}}}-1 \right)+\left( \frac{{{v}^{2}}}{{{c}^{2}}}-\frac{\overrightarrow{a}\cdot \overrightarrow{R}}{{{c}^{2}}}-1 \right)\frac{\partial }{\partial x}\left( \frac{R}{{{K}^{3}}}\frac{\overrightarrow{v}}{c} \right)+\frac{\partial }{\partial x}\left( \frac{1}{{{K}^{2}}}\frac{\overrightarrow{v}}{c} \right)-\frac{\partial }{\partial x}\left( \frac{R}{{{K}^{2}}}\frac{\overrightarrow{a}}{{{c}^{2}}} \right)\]
\[\frac{\partial \overrightarrow{v}}{\partial x}=\frac{\partial \overrightarrow{v}}{\partial t'}\frac{\partial t'}{\partial x}=\overrightarrow{a}\left( -\frac{1}{c}\frac{\partial R}{\partial x} \right)=-\frac{\overrightarrow{a}}{c}\frac{{{R}_{x}}}{K}\]
\[\frac{\partial K}{\partial x}=\frac{{{R}_{x}}}{K}\left( 1+\frac{\overrightarrow{R}\cdot \overrightarrow{a}-{{v}^{2}}}{{{c}^{2}}} \right)-\frac{{{v}_{x}}}{c}\]
\[\frac{\partial }{\partial x}\left( \frac{1}{{{K}^{2}}}\frac{\overrightarrow{v}}{c} \right)=-\frac{{{R}_{x}}}{{{K}^{3}}}\frac{\overrightarrow{a}}{{{c}^{2}}}-\frac{2}{{{K}^{3}}}\frac{\overrightarrow{v}}{c}\left\{ \frac{{{R}_{x}}}{K}\left( 1+\frac{\overrightarrow{R}\cdot \overrightarrow{a}-{{v}^{2}}}{{{c}^{2}}} \right)-\frac{{{v}_{x}}}{c} \right\}\]
\[\frac{\partial }{\partial x}\left( 1+\frac{\overrightarrow{a}\cdot \overrightarrow{R}}{{{c}^{2}}}-\frac{\overrightarrow{v}\cdot \overrightarrow{v}}{{{c}^{2}}} \right)=\frac{1}{{{c}^{2}}}\left\{ -\frac{\overrightarrow{R}}{c}\frac{\partial \overrightarrow{a}}{\partial t'}\frac{{{R}_{x}}}{K}+{{a}_{x}}+3\overrightarrow{v}\frac{\overrightarrow{a}}{c}\frac{{{R}_{x}}}{K} \right\}\]
\[\frac{\partial R}{\partial x}=\frac{{{R}_{x}}}{K}\]
\[\frac{\partial }{\partial x}\left( \frac{R}{{{K}^{2}}}\frac{\overrightarrow{a}}{{{c}^{2}}} \right)=\frac{\overrightarrow{a}}{{{c}^{2}}}\frac{\partial }{\partial x}\left( \frac{R}{{{K}^{2}}} \right)+\frac{R}{{{c}^{2}}{{K}^{2}}}\frac{\partial \overrightarrow{a}}{\partial x}=\frac{\overrightarrow{a}}{{{c}^{2}}}\frac{\partial }{\partial x}\left( \frac{R}{{{K}^{2}}} \right)+\frac{R}{{{c}^{2}}{{K}^{2}}}\frac{\partial \overrightarrow{a}}{\partial t'}\left( -\frac{1}{c}\frac{\partial R}{\partial x} \right)\]
\[\frac{\partial }{\partial x}\left( \frac{R}{{{K}^{2}}}\frac{\overrightarrow{a}}{{{c}^{2}}} \right)=\frac{\overrightarrow{a}}{{{c}^{2}}}\frac{\partial }{\partial x}\left( \frac{R}{{{K}^{2}}} \right)-\frac{R}{{{c}^{3}}{{K}^{2}}}\frac{\partial \overrightarrow{a}}{\partial t'}\frac{{{R}_{x}}}{K}\]
\[\frac{\partial }{\partial x}\left( \frac{R}{{{K}^{2}}} \right)=\frac{1}{{{K}^{2}}}\frac{\partial R}{\partial x}+R\frac{\partial }{\partial x}\left( \frac{1}{{{K}^{2}}} \right)=\frac{1}{{{K}^{2}}}\frac{{{R}_{x}}}{K}-\frac{2R}{{{K}^{3}}}\left\{ \frac{{{R}_{x}}}{K}\left( 1+\frac{\overrightarrow{R}\cdot \overrightarrow{a}-{{v}^{2}}}{{{c}^{2}}} \right)-\frac{{{v}_{x}}}{c} \right\}\]
\[\frac{\partial }{\partial x}\left( \frac{R}{{{K}^{2}}}\frac{\overrightarrow{a}}{{{c}^{2}}} \right)=\frac{\overrightarrow{a}}{{{c}^{2}}}\left[ \frac{1}{{{K}^{2}}}\frac{{{R}_{x}}}{K}-\frac{2R}{{{K}^{3}}}\left\{ \frac{{{R}_{x}}}{K}\left( 1+\frac{\overrightarrow{R}\cdot \overrightarrow{a}-{{v}^{2}}}{{{c}^{2}}} \right)-\frac{{{v}_{x}}}{c} \right\} \right]-\frac{R}{{{c}^{3}}{{K}^{2}}}\frac{\partial \overrightarrow{a}}{\partial t'}\frac{{{R}_{x}}}{K}\]
\[\frac{\partial }{\partial x}\left( \frac{R}{{{K}^{2}}}\frac{\overrightarrow{a}}{{{c}^{2}}} \right)=\frac{1}{{{K}^{3}}{{c}^{2}}}\left( \overrightarrow{a}\left[ {{R}_{x}}-2R\left\{ \frac{{{R}_{x}}}{K}\left( 1+\frac{\overrightarrow{R}\cdot \overrightarrow{a}-{{v}^{2}}}{{{c}^{2}}} \right)-\frac{{{v}_{x}}}{c} \right\} \right]-\frac{R{{R}_{x}}}{c}\frac{\partial \overrightarrow{a}}{\partial t'} \right)\]
\(\frac{}{{{}^{2}}}-\frac{}{/}\frac{}{{{}^{3}}}=\frac{{{}^{2}}}{{{}^{2}}}\)
\[\frac{\partial }{\partial x}\left( \frac{R}{{{K}^{3}}} \right)=\frac{1}{{{K}^{3}}}\frac{\partial R}{\partial x}+R\frac{\partial }{\partial x}\left( \frac{1}{{{K}^{3}}} \right)=\frac{1}{{{K}^{3}}}\frac{{{R}_{x}}}{K}-\frac{3R}{{{K}^{4}}}\left\{ \frac{{{R}_{x}}}{K}\left( 1+\frac{\overrightarrow{R}\cdot \overrightarrow{a}-{{v}^{2}}}{{{c}^{2}}} \right)-\frac{{{v}_{x}}}{c} \right\}\]
\[\frac{\partial \overrightarrow{v}}{\partial x}=\frac{\partial \overrightarrow{v}}{\partial t'}\frac{\partial t'}{\partial x}=\overrightarrow{a}\left( -\frac{1}{c}\frac{\partial R}{\partial x} \right)=-\frac{\overrightarrow{a}}{c}\frac{{{R}_{x}}}{K}\]
\[\frac{\partial }{\partial x}\left( \frac{R}{{{K}^{3}}}\frac{\overrightarrow{v}}{c} \right)=\frac{\overrightarrow{v}}{c}\frac{\partial }{\partial x}\left( \frac{R}{{{K}^{3}}} \right)+\frac{R}{c{{K}^{3}}}\frac{\partial \overrightarrow{v}}{\partial t'}\frac{\partial t'}{\partial x}=\frac{\overrightarrow{v}}{c}\frac{\partial }{\partial x}\left( \frac{R}{{{K}^{3}}} \right)-\frac{R}{c{{K}^{3}}}\frac{\overrightarrow{a}}{c}\frac{{{R}_{x}}}{K}\]
\[\frac{\partial }{\partial x}\left( \frac{R}{{{K}^{3}}}\frac{\overrightarrow{v}}{c} \right)=\frac{\overrightarrow{v}}{c}\left[ \frac{{{R}_{x}}}{{{K}^{4}}}-\frac{3R}{{{K}^{4}}}\left\{ \frac{{{R}_{x}}}{K}\left( 1+\frac{\overrightarrow{R}\cdot \overrightarrow{a}-{{v}^{2}}}{{{c}^{2}}} \right)-\frac{{{v}_{x}}}{c} \right\} \right]-\frac{R{{R}_{x}}}{c{{K}^{4}}}\frac{\overrightarrow{a}}{c}\]
\[\frac{\partial }{\partial x}\left( \frac{R}{{{K}^{3}}}\frac{\overrightarrow{v}}{c} \right)=\frac{1}{{{K}^{4}}c}\left( \overrightarrow{v}\left[ {{R}_{x}}-3R\left\{ \frac{{{R}_{x}}}{K}\left( 1+\frac{\overrightarrow{R}\cdot \overrightarrow{a}-{{v}^{2}}}{{{c}^{2}}} \right)-\frac{{{v}_{x}}}{c} \right\} \right]-\frac{R{{R}_{x}}}{c}\overrightarrow{a} \right)\]
\[\frac{1}{e}\frac{\partial \overrightarrow{{{E}_{2}}}}{\partial x}=\left( \frac{R}{{{K}^{3}}}\frac{\overrightarrow{v}}{c} \right)\frac{\partial }{\partial x}\left( \frac{{{v}^{2}}}{{{c}^{2}}}-\frac{\overrightarrow{a}\cdot \overrightarrow{R}}{{{c}^{2}}}-1 \right)+\left( \frac{{{v}^{2}}}{{{c}^{2}}}-\frac{\overrightarrow{a}\cdot \overrightarrow{R}}{{{c}^{2}}}-1 \right)\frac{\partial }{\partial x}\left( \frac{R}{{{K}^{3}}}\frac{\overrightarrow{v}}{c} \right)+\frac{\partial }{\partial x}\left( \frac{1}{{{K}^{2}}}\frac{\overrightarrow{v}}{c} \right)-\frac{\partial }{\partial x}\left( \frac{R}{{{K}^{2}}}\frac{\overrightarrow{a}}{{{c}^{2}}} \right)\]

\[\begin{align}
  & \frac{1}{e}\frac{\partial \overrightarrow{{{E}_{2}}}}{\partial x}=\left( \frac{R}{{{K}^{3}}}\frac{\overrightarrow{v}}{c} \right)\frac{1}{{{c}^{2}}}\left\{ \frac{\overrightarrow{R}}{c}\frac{\partial \overrightarrow{a}}{\partial t'}\frac{{{R}_{x}}}{K}-{{a}_{x}}-3\overrightarrow{v}\frac{\overrightarrow{a}}{c}\frac{{{R}_{x}}}{K} \right\}+ \\ 
 & +\left( \frac{{{v}^{2}}}{{{c}^{2}}}-\frac{\overrightarrow{a}\cdot \overrightarrow{R}}{{{c}^{2}}}-1 \right)\frac{1}{{{K}^{4}}c}\left( \overrightarrow{v}\left[ {{R}_{x}}-3R\left\{ \frac{{{R}_{x}}}{K}\left( 1+\frac{\overrightarrow{R}\cdot \overrightarrow{a}-{{v}^{2}}}{{{c}^{2}}} \right)-\frac{{{v}_{x}}}{c} \right\} \right]-\frac{R{{R}_{x}}}{c}\overrightarrow{a} \right)- \\ 
 & -\frac{{{R}_{x}}}{{{K}^{3}}}\frac{\overrightarrow{a}}{{{c}^{2}}}-\frac{2}{{{K}^{3}}}\frac{\overrightarrow{v}}{c}\left\{ \frac{{{R}_{x}}}{K}\left( 1+\frac{\overrightarrow{R}\cdot \overrightarrow{a}-{{v}^{2}}}{{{c}^{2}}} \right)-\frac{{{v}_{x}}}{c} \right\}- \\ 
 & -\frac{1}{{{K}^{3}}{{c}^{2}}}\left( \overrightarrow{a}\left[ {{R}_{x}}-2R\left\{ \frac{{{R}_{x}}}{K}\left( 1+\frac{\overrightarrow{R}\cdot \overrightarrow{a}-{{v}^{2}}}{{{c}^{2}}} \right)-\frac{{{v}_{x}}}{c} \right\} \right]-\frac{R{{R}_{x}}}{c}\frac{\partial \overrightarrow{a}}{\partial t'} \right) \\ 
\end{align}\]

\[\begin{align}
  & \frac{1}{e}\frac{\partial \overrightarrow{{{E}_{2}}}}{\partial x}=\frac{R}{{{K}^{3}}}\frac{1}{{{c}^{2}}}\left\{ \left( \frac{\overrightarrow{R}}{c}\frac{\partial \overrightarrow{a}}{\partial t'} \right)\frac{{{R}_{x}}}{K}\frac{\overrightarrow{v}}{c}-{{a}_{x}}\frac{\overrightarrow{v}}{c}-3\left( \overrightarrow{v}\frac{\overrightarrow{a}}{c} \right)\frac{{{R}_{x}}}{K}\frac{\overrightarrow{v}}{c} \right\}- \\ 
 & -\left( 1+\frac{\overrightarrow{R}\cdot \overrightarrow{a}-{{v}^{2}}}{{{c}^{2}}} \right)\frac{1}{{{K}^{4}}}\left( {{R}_{x}}\frac{\overrightarrow{v}}{c}-3R\left\{ \frac{{{R}_{x}}}{K}\frac{\overrightarrow{v}}{c}\left( 1+\frac{\overrightarrow{R}\cdot \overrightarrow{a}-{{v}^{2}}}{{{c}^{2}}} \right)-\frac{{{v}_{x}}}{c}\frac{\overrightarrow{v}}{c} \right\}-\frac{R{{R}_{x}}}{{{c}^{2}}}\overrightarrow{a} \right)- \\ 
 & -\frac{{{R}_{x}}}{{{K}^{3}}}\frac{\overrightarrow{a}}{{{c}^{2}}}-\frac{2}{{{K}^{3}}}\left\{ \frac{{{R}_{x}}}{K}\frac{\overrightarrow{v}}{c}\left( 1+\frac{\overrightarrow{R}\cdot \overrightarrow{a}-{{v}^{2}}}{{{c}^{2}}} \right)-\frac{{{v}_{x}}}{c}\frac{\overrightarrow{v}}{c} \right\}- \\ 
 & -\frac{1}{{{K}^{3}}{{c}^{2}}}\left( {{R}_{x}}\overrightarrow{a}-2R\left\{ \frac{{{R}_{x}}}{K}\overrightarrow{a}\left( 1+\frac{\overrightarrow{R}\cdot \overrightarrow{a}-{{v}^{2}}}{{{c}^{2}}} \right)-\frac{{{v}_{x}}}{c}\overrightarrow{a} \right\}-\frac{R{{R}_{x}}}{c}\frac{\partial \overrightarrow{a}}{\partial t'} \right) \\ 
\end{align}\]

\[\begin{align}
  & \frac{1}{e}\frac{\partial \overrightarrow{{{E}_{2}}}}{\partial x}=\frac{R}{{{K}^{3}}}\frac{1}{{{c}^{2}}}\left\{ \left( \frac{\overrightarrow{R}}{c}\frac{\partial \overrightarrow{a}}{\partial t'} \right)\frac{{{R}_{x}}}{K}\frac{\overrightarrow{v}}{c}-{{a}_{x}}\frac{\overrightarrow{v}}{c}-3\left( \overrightarrow{v}\frac{\overrightarrow{a}}{c} \right)\frac{{{R}_{x}}}{K}\frac{\overrightarrow{v}}{c} \right\}- \\ 
 & -\left( 1+\frac{\overrightarrow{R}\cdot \overrightarrow{a}-{{v}^{2}}}{{{c}^{2}}} \right)\frac{1}{{{K}^{4}}}\left( {{R}_{x}}\frac{\overrightarrow{v}}{c}-3R\left\{ \frac{{{R}_{x}}}{K}\frac{\overrightarrow{v}}{c}\left( 1+\frac{\overrightarrow{R}\cdot \overrightarrow{a}-{{v}^{2}}}{{{c}^{2}}} \right)-\frac{{{v}_{x}}}{c}\frac{\overrightarrow{v}}{c} \right\}-\frac{R{{R}_{x}}}{{{c}^{2}}}\overrightarrow{a} \right)- \\ 
 & -\frac{{{R}_{x}}}{{{K}^{3}}}\frac{\overrightarrow{a}}{{{c}^{2}}}-\frac{2}{{{K}^{3}}}\left\{ \frac{{{R}_{x}}}{K}\frac{\overrightarrow{v}}{c}\left( 1+\frac{\overrightarrow{R}\cdot \overrightarrow{a}-{{v}^{2}}}{{{c}^{2}}} \right)-\frac{{{v}_{x}}}{c}\frac{\overrightarrow{v}}{c} \right\}- \\ 
 & -\frac{{{R}_{x}}\overrightarrow{a}}{{{K}^{3}}{{c}^{2}}}+\frac{R}{{{K}^{3}}{{c}^{2}}}\left\{ 2\frac{{{R}_{x}}}{K}\overrightarrow{a}\left( 1+\frac{\overrightarrow{R}\cdot \overrightarrow{a}-{{v}^{2}}}{{{c}^{2}}} \right)-2\frac{{{v}_{x}}}{c}\overrightarrow{a} \right\}+\frac{R}{{{K}^{3}}{{c}^{2}}}\frac{{{R}_{x}}}{c}\frac{\partial \overrightarrow{a}}{\partial t'} \\ 
\end{align}\]

\[\begin{align}
  & \frac{1}{e}\frac{\partial \overrightarrow{{{E}_{2}}}}{\partial x}=\frac{R}{{{K}^{3}}}\frac{1}{{{c}^{2}}}\left\{ \left( \frac{\overrightarrow{R}}{c}\frac{\partial \overrightarrow{a}}{\partial t'} \right)\frac{{{R}_{x}}}{K}\frac{\overrightarrow{v}}{c}+\frac{{{R}_{x}}}{c}\frac{\partial \overrightarrow{a}}{\partial t'}-{{a}_{x}}\frac{\overrightarrow{v}}{c}-2\frac{{{v}_{x}}}{c}\overrightarrow{a}-3\left( \overrightarrow{v}\frac{\overrightarrow{a}}{c} \right)\frac{{{R}_{x}}}{K}\frac{\overrightarrow{v}}{c} \right\}- \\ 
 & -\left( 1+\frac{\overrightarrow{R}\cdot \overrightarrow{a}-{{v}^{2}}}{{{c}^{2}}} \right)\frac{1}{{{K}^{4}}}\left( {{R}_{x}}\frac{\overrightarrow{v}}{c}-3R\left\{ \frac{{{R}_{x}}}{K}\frac{\overrightarrow{v}}{c}\left( 1+\frac{\overrightarrow{R}\cdot \overrightarrow{a}-{{v}^{2}}}{{{c}^{2}}} \right)-\frac{{{v}_{x}}}{c}\frac{\overrightarrow{v}}{c} \right\}-\frac{R{{R}_{x}}}{{{c}^{2}}}\overrightarrow{a} \right)- \\ 
 & -2\frac{{{R}_{x}}}{{{K}^{3}}}\frac{\overrightarrow{a}}{{{c}^{2}}}-\frac{2}{{{K}^{3}}}\left\{ \frac{{{R}_{x}}}{K}\frac{\overrightarrow{v}}{c}\left( 1+\frac{\overrightarrow{R}\cdot \overrightarrow{a}-{{v}^{2}}}{{{c}^{2}}} \right)-\frac{{{v}_{x}}}{c}\frac{\overrightarrow{v}}{c} \right\}+\frac{2R}{{{K}^{3}}{{c}^{2}}}\frac{{{R}_{x}}}{K}\overrightarrow{a}\left( 1+\frac{\overrightarrow{R}\cdot \overrightarrow{a}-{{v}^{2}}}{{{c}^{2}}} \right) \\ 
\end{align}\]

\[\begin{align}
  & \frac{1}{e}\frac{\partial \overrightarrow{{{E}_{2}}}}{\partial x}=\frac{R}{{{K}^{3}}}\frac{1}{{{c}^{2}}}\left\{ \left( \frac{\overrightarrow{R}}{c}\frac{\partial \overrightarrow{a}}{\partial t'} \right)\frac{{{R}_{x}}}{K}\frac{\overrightarrow{v}}{c}+\frac{{{R}_{x}}}{c}\frac{\partial \overrightarrow{a}}{\partial t'}-{{a}_{x}}\frac{\overrightarrow{v}}{c}-2\frac{{{v}_{x}}}{c}\overrightarrow{a}-3\left( \overrightarrow{v}\frac{\overrightarrow{a}}{c} \right)\frac{{{R}_{x}}}{K}\frac{\overrightarrow{v}}{c} \right\}- \\ 
 & -\left( 1+\frac{\overrightarrow{R}\cdot \overrightarrow{a}-{{v}^{2}}}{{{c}^{2}}} \right)\left( \frac{{{R}_{x}}}{{{K}^{4}}}\frac{\overrightarrow{v}}{c}-\frac{3R}{{{K}^{4}}}\left\{ \frac{{{R}_{x}}}{K}\frac{\overrightarrow{v}}{c}\left( 1+\frac{\overrightarrow{R}\cdot \overrightarrow{a}-{{v}^{2}}}{{{c}^{2}}} \right)-\frac{{{v}_{x}}}{c}\frac{\overrightarrow{v}}{c} \right\}-\frac{R{{R}_{x}}}{{{K}^{4}}{{c}^{2}}}\overrightarrow{a}-\frac{2R}{{{K}^{3}}{{c}^{2}}}\frac{{{R}_{x}}}{K}\overrightarrow{a} \right)- \\ 
 & -2\frac{{{R}_{x}}}{{{K}^{3}}}\frac{\overrightarrow{a}}{{{c}^{2}}}-2\frac{{{R}_{x}}}{{{K}^{4}}}\frac{\overrightarrow{v}}{c}\left( 1+\frac{\overrightarrow{R}\cdot \overrightarrow{a}-{{v}^{2}}}{{{c}^{2}}} \right)+\frac{2}{{{K}^{3}}}\frac{{{v}_{x}}}{c}\frac{\overrightarrow{v}}{c} \\ 
\end{align}\]

\[\begin{align}
  & \frac{1}{e}\frac{\partial \overrightarrow{{{E}_{2}}}}{\partial x}=\frac{R}{{{K}^{3}}}\frac{1}{{{c}^{2}}}\left\{ \left( \frac{\overrightarrow{R}}{c}\frac{\partial \overrightarrow{a}}{\partial t'} \right)\frac{{{R}_{x}}}{K}\frac{\overrightarrow{v}}{c}+\frac{{{R}_{x}}}{c}\frac{\partial \overrightarrow{a}}{\partial t'}-{{a}_{x}}\frac{\overrightarrow{v}}{c}-2\frac{{{v}_{x}}}{c}\overrightarrow{a}-3\left( \overrightarrow{v}\frac{\overrightarrow{a}}{c} \right)\frac{{{R}_{x}}}{K}\frac{\overrightarrow{v}}{c} \right\}- \\ 
 & -\left( 1+\frac{\overrightarrow{R}\cdot \overrightarrow{a}-{{v}^{2}}}{{{c}^{2}}} \right)\left( 3\frac{{{R}_{x}}}{{{K}^{4}}}\frac{\overrightarrow{v}}{c}-\frac{3R}{{{K}^{4}}}\left\{ \frac{{{R}_{x}}}{K}\frac{\overrightarrow{v}}{c}\left( 1+\frac{\overrightarrow{R}\cdot \overrightarrow{a}-{{v}^{2}}}{{{c}^{2}}} \right)-\frac{{{v}_{x}}}{c}\frac{\overrightarrow{v}}{c} \right\}-\frac{3R}{{{K}^{3}}{{c}^{2}}}\frac{{{R}_{x}}}{K}\overrightarrow{a} \right)-2\frac{{{R}_{x}}}{{{K}^{3}}}\frac{\overrightarrow{a}}{{{c}^{2}}}+\frac{2}{{{K}^{3}}}\frac{{{v}_{x}}}{c}\frac{\overrightarrow{v}}{c} \\ 
\end{align}\] Выражения производной второго слагаемого электрического
поля по остальным координатам \(y\) и \(z\) точки наблюдения легко
получить меняя в данном выражении \({{R}_{x}}\) \({{a}_{x}}\)
\({{v}_{x}}\) \(\overrightarrow{i}\) на \({{R}_{y}}\) \({{a}_{y}}\)
\({{v}_{y}}\) \(\overrightarrow{j}\)и на \({{R}_{z}}\) \({{a}_{z}}\)
\({{v}_{z}}\) \(\overrightarrow{k}\)соответственно Возьмём теперь
производную \(x\) компоненты второго слагаемого электрического поля по
\(x\) координате точки наблюдения \[\begin{align}
  & \frac{1}{e}{{\left( \frac{\partial {{E}_{x}}}{\partial x} \right)}_{2}}=\frac{R}{{{K}^{3}}}\frac{1}{{{c}^{2}}}\left\{ \left( \frac{\overrightarrow{R}}{c}\frac{\partial \overrightarrow{a}}{\partial t'} \right)\frac{{{R}_{x}}}{K}\frac{{{v}_{x}}}{c}+\frac{{{R}_{x}}}{c}\frac{\partial {{a}_{x}}}{\partial t'}-{{a}_{x}}\frac{{{v}_{x}}}{c}-2\frac{{{v}_{x}}}{c}{{a}_{x}}-3\left( \overrightarrow{v}\frac{\overrightarrow{a}}{c} \right)\frac{{{R}_{x}}}{K}\frac{{{v}_{x}}}{c} \right\}- \\ 
 & -\left( 1+\frac{\overrightarrow{R}\cdot \overrightarrow{a}-{{v}^{2}}}{{{c}^{2}}} \right)\left( 3\frac{{{R}_{x}}}{{{K}^{4}}}\frac{{{v}_{x}}}{c}-\frac{3R}{{{K}^{4}}}\left\{ \frac{{{R}_{x}}}{K}\frac{{{v}_{x}}}{c}\left( 1+\frac{\overrightarrow{R}\cdot \overrightarrow{a}-{{v}^{2}}}{{{c}^{2}}} \right)-\frac{{{v}_{x}}}{c}\frac{{{v}_{x}}}{c} \right\}-\frac{3R}{{{K}^{3}}{{c}^{2}}}\frac{{{R}_{x}}}{K}{{a}_{x}} \right)-2\frac{{{R}_{x}}}{{{K}^{3}}}\frac{{{a}_{x}}}{{{c}^{2}}}+\frac{2}{{{K}^{3}}}\frac{{{v}_{x}}}{c}\frac{{{x}_{x}}}{c} \\ 
\end{align}\]

\[\begin{align}
  & \frac{1}{e}\frac{\partial {{E}_{2}}_{_{x}}}{\partial x}=\frac{R}{{{K}^{3}}}\frac{1}{{{c}^{2}}}\left\{ \left( \frac{\overrightarrow{R}}{c}\frac{\partial \overrightarrow{a}}{\partial t'} \right)\frac{{{R}_{x}}}{K}\frac{{{v}_{x}}}{c}+\frac{{{R}_{x}}}{c}\frac{\partial {{a}_{x}}}{\partial t'}-3\frac{{{v}_{x}}}{c}{{a}_{x}}-3\left( \overrightarrow{v}\frac{\overrightarrow{a}}{c} \right)\frac{{{R}_{x}}}{K}\frac{{{v}_{x}}}{c} \right\}- \\ 
 & -\left( 1+\frac{\overrightarrow{R}\cdot \overrightarrow{a}-{{v}^{2}}}{{{c}^{2}}} \right)\left( 3\frac{{{R}_{x}}}{{{K}^{4}}}\frac{{{v}_{x}}}{c}-\frac{3R}{{{K}^{4}}}\left\{ \frac{{{R}_{x}}}{K}\frac{{{v}_{x}}}{c}\left( 1+\frac{\overrightarrow{R}\cdot \overrightarrow{a}-{{v}^{2}}}{{{c}^{2}}} \right)-\frac{{{v}_{x}}}{c}\frac{{{v}_{x}}}{c} \right\}-\frac{3R}{{{K}^{3}}{{c}^{2}}}\frac{{{R}_{x}}}{K}{{a}_{x}} \right)-2\frac{{{R}_{x}}}{{{K}^{3}}}\frac{{{a}_{x}}}{{{c}^{2}}}+\frac{2}{{{K}^{3}}}\frac{{{v}_{x}}}{c}\frac{{{x}_{x}}}{c} \\ 
\end{align}\]

\[\begin{align}
  & \frac{1}{e}\frac{\partial {{E}_{2}}_{_{x}}}{\partial x}=\frac{R}{{{K}^{3}}}\frac{1}{{{c}^{2}}}\left\{ \left( \frac{\overrightarrow{R}}{c}\frac{\partial \overrightarrow{a}}{\partial t'} \right)\frac{{{R}_{x}}}{K}\frac{{{v}_{x}}}{c}+\frac{{{R}_{x}}}{c}\frac{\partial {{a}_{x}}}{\partial t'}-3\frac{{{v}_{x}}}{c}{{a}_{x}}-3\left( \overrightarrow{v}\frac{\overrightarrow{a}}{c} \right)\frac{{{R}_{x}}}{K}\frac{{{v}_{x}}}{c} \right\}- \\ 
 & -\left( 1+\frac{\overrightarrow{R}\cdot \overrightarrow{a}-{{v}^{2}}}{{{c}^{2}}} \right)\frac{1}{{{K}^{3}}}\left( 3\frac{{{R}_{x}}}{K}\frac{{{v}_{x}}}{c}-\frac{3R}{K}\left\{ \frac{{{R}_{x}}}{K}\frac{{{v}_{x}}}{c}\left( 1+\frac{\overrightarrow{R}\cdot \overrightarrow{a}-{{v}^{2}}}{{{c}^{2}}} \right)-\frac{{{v}_{x}}}{c}\frac{{{v}_{x}}}{c} \right\}-\frac{3R}{K}\frac{{{R}_{x}}}{{{c}^{2}}}{{a}_{x}} \right)-2\frac{{{R}_{x}}}{{{K}^{3}}}\frac{{{a}_{x}}}{{{c}^{2}}}+\frac{2}{{{K}^{3}}}\frac{{{v}_{x}}}{c}\frac{{{x}_{x}}}{c} \\ 
\end{align}\] Выражения производных \(y\) и \(z\) компонент второго
слагаемого электрического поля по \(y\)и \(z\) координатам точки
наблюдения легко получить аналогичной заменой индексов в данном
выражении Скалярно складывая производные\(x\) , \(y\) и \(z\)компонент
второго слагаемого электрического поля по соответственно \(x\) , \(y\) и
\(z\)координатам точки наблюдения получаем выражение для дивергенции
второго слагаемого электрического поля \[\begin{align}
  & \frac{1}{e}{{\left( \frac{\partial {{E}_{x}}}{\partial x}+\frac{\partial {{E}_{y}}}{\partial y}+\frac{\partial {{E}_{z}}}{\partial z} \right)}_{2}}=\frac{R}{{{K}^{3}}}\frac{1}{{{c}^{2}}}\left\{ \left( \frac{\overrightarrow{R}}{c}\frac{\partial \overrightarrow{a}}{\partial t'} \right)\frac{\overrightarrow{R}}{K}\frac{\overrightarrow{v}}{c}+\frac{\overrightarrow{R}}{c}\frac{\partial \overrightarrow{a}}{\partial t'}-3\frac{\overrightarrow{v}}{c}\overrightarrow{a}-3\left( \overrightarrow{v}\frac{\overrightarrow{a}}{c} \right)\frac{\overrightarrow{R}}{K}\frac{\overrightarrow{v}}{c} \right\}- \\ 
 & -\left( 1+\frac{\overrightarrow{R}\cdot \overrightarrow{a}-{{v}^{2}}}{{{c}^{2}}} \right)\left( 3\frac{\overrightarrow{R}}{{{K}^{4}}}\frac{\overrightarrow{v}}{c}-\frac{3R}{{{K}^{4}}}\left\{ \frac{\overrightarrow{R}}{K}\frac{\overrightarrow{v}}{c}\left( 1+\frac{\overrightarrow{R}\cdot \overrightarrow{a}-{{v}^{2}}}{{{c}^{2}}} \right)-\frac{\overrightarrow{v}}{c}\frac{\overrightarrow{v}}{c} \right\}-\frac{3R\overrightarrow{R}}{{{K}^{4}}{{c}^{2}}}\overrightarrow{a} \right)-2\frac{\overrightarrow{R}}{{{K}^{3}}}\frac{\overrightarrow{a}}{{{c}^{2}}}+\frac{2}{{{K}^{3}}}\frac{\overrightarrow{v}}{c}\frac{\overrightarrow{v}}{c} \\ 
\end{align}\]

\[\begin{align}
  & \frac{1}{e}{{\left( \frac{\partial {{E}_{x}}}{\partial x}+\frac{\partial {{E}_{y}}}{\partial y}+\frac{\partial {{E}_{z}}}{\partial z} \right)}_{2}}=\frac{R}{{{K}^{3}}}\frac{1}{{{c}^{2}}}\left\{ \left( \frac{\overrightarrow{R}}{c}\frac{\partial \overrightarrow{a}}{\partial t'} \right)\left( \frac{\overrightarrow{R}}{K}\frac{\overrightarrow{v}}{c}+1 \right)-3\left( \overrightarrow{v}\frac{\overrightarrow{a}}{c} \right)\left( \frac{\overrightarrow{R}}{K}\frac{\overrightarrow{v}}{c}+1 \right) \right\}- \\ 
 & -\left( 1+\frac{\overrightarrow{R}\cdot \overrightarrow{a}-{{v}^{2}}}{{{c}^{2}}} \right)\left( 3\frac{\overrightarrow{R}}{{{K}^{4}}}\frac{\overrightarrow{v}}{c}-\frac{3R}{{{K}^{4}}}\left\{ \frac{\overrightarrow{R}}{K}\frac{\overrightarrow{v}}{c}\left( 1+\frac{\overrightarrow{R}\cdot \overrightarrow{a}-{{v}^{2}}}{{{c}^{2}}} \right)-\frac{\overrightarrow{v}}{c}\frac{\overrightarrow{v}}{c} \right\}-\frac{3R\overrightarrow{R}}{{{K}^{4}}{{c}^{2}}}\overrightarrow{a} \right)-2\frac{\overrightarrow{R}}{{{K}^{3}}}\frac{\overrightarrow{a}}{{{c}^{2}}}+\frac{2}{{{K}^{3}}}\frac{\overrightarrow{v}}{c}\frac{\overrightarrow{v}}{c} \\ 
\end{align}\]

\[\begin{align}
  & \frac{1}{e}{{\left( \frac{\partial {{E}_{x}}}{\partial x}+\frac{\partial {{E}_{y}}}{\partial y}+\frac{\partial {{E}_{z}}}{\partial z} \right)}_{2}}=\frac{R}{{{K}^{3}}}\frac{1}{{{c}^{2}}}\left\{ \left( \frac{\overrightarrow{R}}{c}\frac{\partial \overrightarrow{a}}{\partial t'} \right)\left( \frac{\overrightarrow{R}}{K}\frac{\overrightarrow{v}}{c}+1 \right)-3\left( \overrightarrow{v}\frac{\overrightarrow{a}}{c} \right)\left( \frac{\overrightarrow{R}}{K}\frac{\overrightarrow{v}}{c}+1 \right) \right\}- \\ 
 & -\left( 1+\frac{\overrightarrow{R}\cdot \overrightarrow{a}-{{v}^{2}}}{{{c}^{2}}} \right)\left( 3\frac{\overrightarrow{R}}{{{K}^{4}}}\frac{\overrightarrow{v}}{c}-\frac{3R}{{{K}^{4}}}\left\{ \frac{\overrightarrow{R}}{K}\frac{\overrightarrow{v}}{c}\left( 1+\frac{\overrightarrow{R}\cdot \overrightarrow{a}-{{v}^{2}}}{{{c}^{2}}} \right)-\frac{\overrightarrow{v}}{c}\frac{\overrightarrow{v}}{c} \right\}-\frac{3R\overrightarrow{R}}{{{K}^{4}}{{c}^{2}}}\overrightarrow{a} \right)-2\frac{\overrightarrow{R}}{{{K}^{3}}}\frac{\overrightarrow{a}}{{{c}^{2}}}+\frac{2}{{{K}^{3}}}\frac{\overrightarrow{v}}{c}\frac{\overrightarrow{v}}{c} \\ 
\end{align}\] Сумма двух компонент электрического поля
\[\overrightarrow{E}=-\nabla \varphi -\frac{1}{c}\frac{\partial \overrightarrow{A}}{\partial t}=\frac{e}{{{K}^{3}}}\left[ \left( \overrightarrow{R}-R\frac{\overrightarrow{v}}{c} \right)\left( 1+\frac{\overrightarrow{a}\cdot \overrightarrow{R}}{{{c}^{2}}}-\frac{\overrightarrow{v}\cdot \overrightarrow{v}}{{{c}^{2}}} \right)-KR\frac{\overrightarrow{a}}{{{c}^{2}}} \right]\]
\[\overrightarrow{E}=\frac{e}{{{K}^{3}}}\left( \overrightarrow{R}-R\frac{\overrightarrow{v}}{c} \right)\left( 1+\frac{\overrightarrow{a}\cdot \overrightarrow{R}}{{{c}^{2}}}-\frac{\overrightarrow{v}\cdot \overrightarrow{v}}{{{c}^{2}}} \right)-\frac{e}{{{K}^{2}}}R\frac{\overrightarrow{a}}{{{c}^{2}}}\]
\[\frac{1}{e}\frac{\partial \overrightarrow{E}}{\partial x}=\frac{\partial }{\partial x}\left[ \frac{1}{{{K}^{3}}}\left( \overrightarrow{R}-R\frac{\overrightarrow{v}}{c} \right)\left( 1+\frac{\overrightarrow{a}\cdot \overrightarrow{R}}{{{c}^{2}}}-\frac{\overrightarrow{v}\cdot \overrightarrow{v}}{{{c}^{2}}} \right)-\frac{R}{{{K}^{2}}}\frac{\overrightarrow{a}}{{{c}^{2}}} \right]\]
\[\frac{1}{e}\frac{\partial \overrightarrow{E}}{\partial x}=\frac{1}{{{K}^{3}}}\left( \overrightarrow{R}-R\frac{\overrightarrow{v}}{c} \right)\frac{\partial }{\partial x}\left( 1+\frac{\overrightarrow{a}\cdot \overrightarrow{R}}{{{c}^{2}}}-\frac{\overrightarrow{v}\cdot \overrightarrow{v}}{{{c}^{2}}} \right)+\left( 1+\frac{\overrightarrow{a}\cdot \overrightarrow{R}}{{{c}^{2}}}-\frac{\overrightarrow{v}\cdot \overrightarrow{v}}{{{c}^{2}}} \right)\frac{\partial }{\partial x}\left\{ \frac{1}{{{K}^{3}}}\left( \overrightarrow{R}-R\frac{\overrightarrow{v}}{c} \right) \right\}-\frac{\partial }{\partial x}\left( \frac{R}{{{K}^{2}}}\frac{\overrightarrow{a}}{{{c}^{2}}} \right)\]
\[\frac{1}{e}\frac{\partial \overrightarrow{E}}{\partial x}=\frac{1}{{{K}^{3}}}\left( \overrightarrow{R}-R\frac{\overrightarrow{v}}{c} \right)\frac{\partial }{\partial x}\left( 1+\frac{\overrightarrow{a}\cdot \overrightarrow{R}}{{{c}^{2}}}-\frac{\overrightarrow{v}\cdot \overrightarrow{v}}{{{c}^{2}}} \right)+\left( 1+\frac{\overrightarrow{a}\cdot \overrightarrow{R}}{{{c}^{2}}}-\frac{\overrightarrow{v}\cdot \overrightarrow{v}}{{{c}^{2}}} \right)\frac{\partial }{\partial x}\left\{ \frac{\overrightarrow{R}}{{{K}^{3}}}-\frac{R}{{{K}^{3}}}\frac{\overrightarrow{v}}{c} \right\}-\frac{\partial }{\partial x}\left( \frac{R}{{{K}^{2}}}\frac{\overrightarrow{a}}{{{c}^{2}}} \right)\]..
\[\frac{\partial }{\partial x}\left( 1+\frac{\overrightarrow{a}\cdot \overrightarrow{R}}{{{c}^{2}}}-\frac{\overrightarrow{v}\cdot \overrightarrow{v}}{{{c}^{2}}} \right)=\frac{1}{{{c}^{2}}}\left\{ -\frac{\overrightarrow{R}}{c}\frac{\partial \overrightarrow{a}}{\partial t'}\frac{{{R}_{x}}}{K}+{{a}_{x}}+3\overrightarrow{v}\frac{\overrightarrow{a}}{c}\frac{{{R}_{x}}}{K} \right\}\].
\[\frac{\partial }{\partial x}\left( \frac{\overrightarrow{R}}{{{K}^{3}}} \right)=\frac{1}{{{K}^{3}}}\left( \overrightarrow{i}+\frac{\overrightarrow{v}}{c}\frac{{{R}_{x}}}{K} \right)-\overrightarrow{R}\frac{3}{{{K}^{4}}}\left( \frac{{{R}_{x}}}{K}+\left( \frac{\overrightarrow{R}\cdot \overrightarrow{a}-{{v}^{2}}}{{{c}^{2}}} \right)\frac{{{R}_{x}}}{K}-\frac{{{v}_{x}}}{c} \right)\]
\[\frac{\partial }{\partial x}\left( \frac{R}{{{K}^{3}}}\frac{\overrightarrow{v}}{c} \right)=\frac{1}{{{K}^{4}}c}\left( \overrightarrow{v}\left[ {{R}_{x}}-3R\left\{ \frac{{{R}_{x}}}{K}\left( 1+\frac{\overrightarrow{R}\cdot \overrightarrow{a}-{{v}^{2}}}{{{c}^{2}}} \right)-\frac{{{v}_{x}}}{c} \right\} \right]-\frac{R{{R}_{x}}}{c}\overrightarrow{a} \right)\]
\[\begin{align}
  & \frac{\partial }{\partial x}\left\{ \frac{\overrightarrow{R}}{{{K}^{3}}}-\frac{R}{{{K}^{3}}}\frac{\overrightarrow{v}}{c} \right\}=\frac{1}{{{K}^{3}}}\left( \overrightarrow{i}+\frac{\overrightarrow{v}}{c}\frac{{{R}_{x}}}{K} \right)-\overrightarrow{R}\frac{3}{{{K}^{4}}}\left( \frac{{{R}_{x}}}{K}+\left( \frac{\overrightarrow{R}\cdot \overrightarrow{a}-{{v}^{2}}}{{{c}^{2}}} \right)\frac{{{R}_{x}}}{K}-\frac{{{v}_{x}}}{c} \right)- \\ 
 & -\frac{1}{{{K}^{4}}c}\left( \overrightarrow{v}\left[ {{R}_{x}}-3R\left\{ \frac{{{R}_{x}}}{K}\left( 1+\frac{\overrightarrow{R}\cdot \overrightarrow{a}-{{v}^{2}}}{{{c}^{2}}} \right)-\frac{{{v}_{x}}}{c} \right\} \right]-\frac{R{{R}_{x}}}{c}\overrightarrow{a} \right)= \\ 
 &  \\ 
 & =\left( \overrightarrow{i}\frac{1}{{{K}^{3}}}+\frac{\overrightarrow{v}}{c}\frac{{{R}_{x}}}{{{K}^{4}}} \right)-\overrightarrow{R}\frac{3}{{{K}^{4}}}\left( \frac{{{R}_{x}}}{K}+\left( \frac{\overrightarrow{R}\cdot \overrightarrow{a}-{{v}^{2}}}{{{c}^{2}}} \right)\frac{{{R}_{x}}}{K}-\frac{{{v}_{x}}}{c} \right)- \\ 
 & -\left[ \frac{{{R}_{x}}}{{{K}^{4}}}\frac{\overrightarrow{v}}{c}-\frac{\overrightarrow{v}}{c}\frac{3R}{{{K}^{4}}}\left\{ \frac{{{R}_{x}}}{K}\left( 1+\frac{\overrightarrow{R}\cdot \overrightarrow{a}-{{v}^{2}}}{{{c}^{2}}} \right)-\frac{{{v}_{x}}}{c} \right\} \right]+\frac{R{{R}_{x}}}{{{K}^{4}}}\frac{\overrightarrow{a}}{{{c}^{2}}}= \\ 
 &  \\ 
 & =\left( \overrightarrow{i}\frac{1}{{{K}^{3}}} \right)-\overrightarrow{R}\frac{3}{{{K}^{4}}}\left( \frac{{{R}_{x}}}{K}\left( 1+\frac{\overrightarrow{R}\cdot \overrightarrow{a}-{{v}^{2}}}{{{c}^{2}}} \right)-\frac{{{v}_{x}}}{c} \right)+ \\ 
 & +\frac{\overrightarrow{v}}{c}\frac{3R}{{{K}^{4}}}\left\{ \frac{{{R}_{x}}}{K}\left( 1+\frac{\overrightarrow{R}\cdot \overrightarrow{a}-{{v}^{2}}}{{{c}^{2}}} \right)-\frac{{{v}_{x}}}{c} \right\}+\frac{R{{R}_{x}}}{{{K}^{4}}}\frac{\overrightarrow{a}}{{{c}^{2}}}= \\ 
 &  \\ 
 & =\left( \overrightarrow{i}\frac{1}{{{K}^{3}}} \right)+\left( \frac{\overrightarrow{v}}{c}\frac{3R}{{{K}^{4}}}-\overrightarrow{R}\frac{3}{{{K}^{4}}} \right)\left\{ \frac{{{R}_{x}}}{K}\left( 1+\frac{\overrightarrow{R}\cdot \overrightarrow{a}-{{v}^{2}}}{{{c}^{2}}} \right)-\frac{{{v}_{x}}}{c} \right\}+\frac{R{{R}_{x}}}{{{K}^{4}}}\frac{\overrightarrow{a}}{{{c}^{2}}} \\ 
\end{align}\]
\[\frac{\partial }{\partial x}\left\{ \frac{\overrightarrow{R}}{{{K}^{3}}}-\frac{R}{{{K}^{3}}}\frac{\overrightarrow{v}}{c} \right\}=\frac{1}{{{K}^{3}}}\left[ \overrightarrow{i}+\left( \frac{\overrightarrow{v}}{c}\frac{3R}{K}-\overrightarrow{R}\frac{3}{K} \right)\left\{ \frac{{{R}_{x}}}{K}\left( 1+\frac{\overrightarrow{R}\cdot \overrightarrow{a}-{{v}^{2}}}{{{c}^{2}}} \right)-\frac{{{v}_{x}}}{c} \right\}+\frac{R{{R}_{x}}}{K}\frac{\overrightarrow{a}}{{{c}^{2}}} \right]\]
\[\frac{\partial }{\partial x}\left( \frac{R}{{{K}^{2}}}\frac{\overrightarrow{a}}{{{c}^{2}}} \right)=\frac{1}{{{K}^{3}}{{c}^{2}}}\left( \overrightarrow{a}\left[ {{R}_{x}}-2R\left\{ \frac{{{R}_{x}}}{K}\left( 1+\frac{\overrightarrow{R}\cdot \overrightarrow{a}-{{v}^{2}}}{{{c}^{2}}} \right)-\frac{{{v}_{x}}}{c} \right\} \right]-\frac{R{{R}_{x}}}{c}\frac{\partial \overrightarrow{a}}{\partial t'} \right)\]
\[\begin{align}
  & \frac{1}{e}\frac{\partial \overrightarrow{E}}{\partial x}=\frac{1}{{{K}^{3}}}\left( \overrightarrow{R}-R\frac{\overrightarrow{v}}{c} \right)\frac{1}{{{c}^{2}}}\left\{ -\left( \frac{\overrightarrow{R}}{c}\frac{\partial \overrightarrow{a}}{\partial t'} \right)\frac{{{R}_{x}}}{K}+{{a}_{x}}+3\left( \overrightarrow{v}\frac{\overrightarrow{a}}{c} \right)\frac{{{R}_{x}}}{K} \right\}+ \\ 
 & +\left( 1+\frac{\overrightarrow{a}\cdot \overrightarrow{R}}{{{c}^{2}}}-\frac{\overrightarrow{v}\cdot \overrightarrow{v}}{{{c}^{2}}} \right)\frac{1}{{{K}^{3}}}\left[ \overrightarrow{i}+\left( \frac{\overrightarrow{v}}{c}\frac{3R}{K}-\overrightarrow{R}\frac{3}{K} \right)\left\{ \frac{{{R}_{x}}}{K}\left( 1+\frac{\overrightarrow{R}\cdot \overrightarrow{a}-{{v}^{2}}}{{{c}^{2}}} \right)-\frac{{{v}_{x}}}{c} \right\}+\frac{R{{R}_{x}}}{K}\frac{\overrightarrow{a}}{{{c}^{2}}} \right]- \\ 
 & -\frac{1}{{{K}^{3}}{{c}^{2}}}\left( \overrightarrow{a}\left[ {{R}_{x}}-2R\left\{ \frac{{{R}_{x}}}{K}\left( 1+\frac{\overrightarrow{R}\cdot \overrightarrow{a}-{{v}^{2}}}{{{c}^{2}}} \right)-\frac{{{v}_{x}}}{c} \right\} \right]-\frac{R{{R}_{x}}}{c}\frac{\partial \overrightarrow{a}}{\partial t'} \right) \\ 
\end{align}\] Выражения производной электрического поля по остальным
координатам \(y\) и \(z\) точки наблюдения легко получить меняя в данном
выражении \({{R}_{x}}\) \({{a}_{x}}\) \({{v}_{x}}\)
\(\overrightarrow{i}\) на \({{R}_{y}}\) \({{a}_{y}}\) \({{v}_{y}}\)
\(\overrightarrow{j}\)и на \({{R}_{z}}\) \({{a}_{z}}\) \({{v}_{z}}\)
\(\overrightarrow{k}\)соответственно Возьмём теперь производную
компоненты электрического поля по координате точки наблюдения
\[\begin{align}
  & \frac{1}{e}\frac{\partial {{E}_{x}}}{\partial x}=\frac{1}{{{K}^{3}}}\frac{1}{{{c}^{2}}}\left( {{R}_{x}}-R\frac{{{v}_{x}}}{c} \right)\left\{ -\left( \frac{\overrightarrow{R}}{c}\frac{\partial \overrightarrow{a}}{\partial t'} \right)\frac{{{R}_{x}}}{K}+{{a}_{x}}+3\left( \overrightarrow{v}\frac{\overrightarrow{a}}{c} \right)\frac{{{R}_{x}}}{K} \right\}+ \\ 
 & +\left( 1+\frac{\overrightarrow{a}\cdot \overrightarrow{R}}{{{c}^{2}}}-\frac{\overrightarrow{v}\cdot \overrightarrow{v}}{{{c}^{2}}} \right)\frac{1}{{{K}^{3}}}\left[ 1+\left[ \left( \frac{{{v}_{x}}}{c}\frac{3R}{K}-{{R}_{x}}\frac{3}{K} \right)\left\{ \frac{{{R}_{x}}}{K}\left( 1+\frac{\overrightarrow{R}\cdot \overrightarrow{a}-{{v}^{2}}}{{{c}^{2}}} \right)-\frac{{{v}_{x}}}{c} \right\} \right]+\frac{R{{R}_{x}}}{K}\frac{{{a}_{x}}}{{{c}^{2}}} \right]- \\ 
 & -\frac{1}{{{K}^{3}}{{c}^{2}}}\left( {{a}_{x}}\left[ {{R}_{x}}-2R\left\{ \frac{{{R}_{x}}}{K}\left( 1+\frac{\overrightarrow{R}\cdot \overrightarrow{a}-{{v}^{2}}}{{{c}^{2}}} \right)-\frac{{{v}_{x}}}{c} \right\} \right]-\frac{R{{R}_{x}}}{c}\frac{\partial {{a}_{x}}}{\partial t'} \right) \\ 
\end{align}\]

\[\begin{align}
  & \frac{1}{e}\frac{\partial {{E}_{x}}}{\partial x}=\frac{1}{{{c}^{2}}}\frac{1}{{{K}^{3}}}\left[ {{R}_{x}}\left\{ -\left( \frac{\overrightarrow{R}}{c}\frac{\partial \overrightarrow{a}}{\partial t'} \right)\frac{{{R}_{x}}}{K}+{{a}_{x}}+3\left( \overrightarrow{v}\frac{\overrightarrow{a}}{c} \right)\frac{{{R}_{x}}}{K} \right\}-R\frac{{{v}_{x}}}{c}\left\{ -\left( \frac{\overrightarrow{R}}{c}\frac{\partial \overrightarrow{a}}{\partial t'} \right)\frac{{{R}_{x}}}{K}+{{a}_{x}}+3\left( \overrightarrow{v}\frac{\overrightarrow{a}}{c} \right)\frac{{{R}_{x}}}{K} \right\} \right]+ \\ 
 & +\left( 1+\frac{\overrightarrow{a}\cdot \overrightarrow{R}}{{{c}^{2}}}-\frac{\overrightarrow{v}\cdot \overrightarrow{v}}{{{c}^{2}}} \right)\frac{1}{{{K}^{3}}}\left[ 1+\left[ \left( \frac{{{v}_{x}}}{c}\frac{3R}{K}-{{R}_{x}}\frac{3}{K} \right)\left\{ \frac{{{R}_{x}}}{K}\left( 1+\frac{\overrightarrow{R}\cdot \overrightarrow{a}-{{v}^{2}}}{{{c}^{2}}} \right)-\frac{{{v}_{x}}}{c} \right\} \right]+\frac{R{{R}_{x}}}{K}\frac{{{a}_{x}}}{{{c}^{2}}} \right]- \\ 
 & -\frac{1}{{{K}^{3}}{{c}^{2}}}\left( {{a}_{x}}{{R}_{x}}-2R\left\{ \frac{{{a}_{x}}{{R}_{x}}}{K}\left( 1+\frac{\overrightarrow{R}\cdot \overrightarrow{a}-{{v}^{2}}}{{{c}^{2}}} \right)-\frac{{{a}_{x}}{{v}_{x}}}{c} \right\}-\frac{R{{R}_{x}}}{c}\frac{\partial {{a}_{x}}}{\partial t'} \right) \\ 
\end{align}\]

\[\begin{align}
  & \frac{1}{e}\frac{\partial {{E}_{x}}}{\partial x}=\frac{1}{{{c}^{2}}}\frac{1}{{{K}^{3}}}\left[ -\left( \frac{\overrightarrow{R}}{c}\frac{\partial \overrightarrow{a}}{\partial t'} \right)\frac{{{R}_{x}}{{R}_{x}}}{K}+{{a}_{x}}{{R}_{x}}+3\left( \overrightarrow{v}\frac{\overrightarrow{a}}{c} \right)\frac{{{R}_{x}}{{R}_{x}}}{K}+\left( \frac{\overrightarrow{R}}{c}\frac{\partial \overrightarrow{a}}{\partial t'} \right)\frac{R}{c}\frac{{{R}_{x}}{{v}_{x}}}{K}-\frac{R}{c}{{a}_{x}}{{v}_{x}}-3\frac{R}{c}\left( \overrightarrow{v}\frac{\overrightarrow{a}}{c} \right)\frac{{{R}_{x}}{{v}_{x}}}{K} \right]+ \\ 
 & +\left( 1+\frac{\overrightarrow{a}\cdot \overrightarrow{R}}{{{c}^{2}}}-\frac{\overrightarrow{v}\cdot \overrightarrow{v}}{{{c}^{2}}} \right)\frac{1}{{{K}^{3}}}\left[ 1+\left[ \left( \frac{{{v}_{x}}}{c}\frac{3R}{K}-{{R}_{x}}\frac{3}{K} \right)\left\{ \frac{{{R}_{x}}}{K}\left( 1+\frac{\overrightarrow{R}\cdot \overrightarrow{a}-{{v}^{2}}}{{{c}^{2}}} \right)-\frac{{{v}_{x}}}{c} \right\} \right]+\frac{R{{R}_{x}}}{K}\frac{{{a}_{x}}}{{{c}^{2}}} \right]- \\ 
 & -\frac{1}{{{K}^{3}}{{c}^{2}}}\left( {{a}_{x}}{{R}_{x}}-2R\left\{ \frac{{{a}_{x}}{{R}_{x}}}{K}\left( 1+\frac{\overrightarrow{R}\cdot \overrightarrow{a}-{{v}^{2}}}{{{c}^{2}}} \right)-\frac{{{a}_{x}}{{v}_{x}}}{c} \right\}-\frac{R{{R}_{x}}}{c}\frac{\partial {{a}_{x}}}{\partial t'} \right) \\ 
\end{align}\]

\[\begin{align}
  & \frac{1}{e}\frac{\partial {{E}_{x}}}{\partial x}=\frac{1}{{{c}^{2}}}\frac{1}{{{K}^{3}}}\left[ -\left( \frac{\overrightarrow{R}}{c}\frac{\partial \overrightarrow{a}}{\partial t'} \right)\frac{{{R}_{x}}{{R}_{x}}}{K}+3\left( \overrightarrow{v}\frac{\overrightarrow{a}}{c} \right)\frac{{{R}_{x}}{{R}_{x}}}{K}+\left( \frac{\overrightarrow{R}}{c}\frac{\partial \overrightarrow{a}}{\partial t'} \right)\frac{R}{c}\frac{{{R}_{x}}{{v}_{x}}}{K}-\frac{R}{c}{{a}_{x}}{{v}_{x}}-3\frac{R}{c}\left( \overrightarrow{v}\frac{\overrightarrow{a}}{c} \right)\frac{{{R}_{x}}{{v}_{x}}}{K} \right]+ \\ 
 & +\left( 1+\frac{\overrightarrow{a}\cdot \overrightarrow{R}}{{{c}^{2}}}-\frac{\overrightarrow{v}\cdot \overrightarrow{v}}{{{c}^{2}}} \right)\frac{1}{{{K}^{3}}}\left[ 1+\left[ \left( \frac{{{v}_{x}}}{c}\frac{3R}{K}-{{R}_{x}}\frac{3}{K} \right)\left\{ \frac{{{R}_{x}}}{K}\left( 1+\frac{\overrightarrow{R}\cdot \overrightarrow{a}-{{v}^{2}}}{{{c}^{2}}} \right)-\frac{{{v}_{x}}}{c} \right\} \right]+\frac{R{{R}_{x}}}{K}\frac{{{a}_{x}}}{{{c}^{2}}} \right]+ \\ 
 & +\frac{1}{{{K}^{3}}{{c}^{2}}}\left( 2R\left\{ \frac{{{a}_{x}}{{R}_{x}}}{K}\left( 1+\frac{\overrightarrow{R}\cdot \overrightarrow{a}-{{v}^{2}}}{{{c}^{2}}} \right)-\frac{{{a}_{x}}{{v}_{x}}}{c} \right\}+\frac{R{{R}_{x}}}{c}\frac{\partial {{a}_{x}}}{\partial t'} \right) \\ 
\end{align}\]

\[\begin{align}
  & \frac{1}{e}\frac{\partial {{E}_{x}}}{\partial x}=\frac{1}{{{c}^{2}}}\frac{1}{{{K}^{3}}}\left[ -\left( \frac{\overrightarrow{R}}{c}\frac{\partial \overrightarrow{a}}{\partial t'} \right)\frac{{{R}_{x}}{{R}_{x}}}{K}+3\left( \overrightarrow{v}\frac{\overrightarrow{a}}{c} \right)\frac{{{R}_{x}}{{R}_{x}}}{K}+\left( \frac{\overrightarrow{R}}{c}\frac{\partial \overrightarrow{a}}{\partial t'} \right)\frac{R}{c}\frac{{{R}_{x}}{{v}_{x}}}{K}-\frac{R}{c}{{a}_{x}}{{v}_{x}}-3\frac{R}{c}\left( \overrightarrow{v}\frac{\overrightarrow{a}}{c} \right)\frac{{{R}_{x}}{{v}_{x}}}{K} \right]+ \\ 
 & +\left( 1+\frac{\overrightarrow{a}\cdot \overrightarrow{R}}{{{c}^{2}}}-\frac{\overrightarrow{v}\cdot \overrightarrow{v}}{{{c}^{2}}} \right)\frac{1}{{{K}^{3}}}\left[ 1+\frac{{{v}_{x}}}{c}\frac{3R}{K}\left\{ \frac{{{R}_{x}}}{K}\left( 1+\frac{\overrightarrow{R}\cdot \overrightarrow{a}-{{v}^{2}}}{{{c}^{2}}} \right)-\frac{{{v}_{x}}}{c} \right\}-{{R}_{x}}\frac{3}{K}\left\{ \frac{{{R}_{x}}}{K}\left( 1+\frac{\overrightarrow{R}\cdot \overrightarrow{a}-{{v}^{2}}}{{{c}^{2}}} \right)-\frac{{{v}_{x}}}{c} \right\}+\frac{R{{R}_{x}}}{K}\frac{{{a}_{x}}}{{{c}^{2}}} \right]+ \\ 
 & +\frac{1}{{{K}^{3}}{{c}^{2}}}\left( 2R\frac{{{a}_{x}}{{R}_{x}}}{K}\left( 1+\frac{\overrightarrow{R}\cdot \overrightarrow{a}-{{v}^{2}}}{{{c}^{2}}} \right)-\frac{2R}{c}{{a}_{x}}{{v}_{x}}+\frac{R{{R}_{x}}}{c}\frac{\partial {{a}_{x}}}{\partial t'} \right) \\ 
\end{align}\]

\[\begin{align}
  & \frac{1}{e}\frac{\partial {{E}_{x}}}{\partial x}=\frac{1}{{{c}^{2}}}\frac{1}{{{K}^{3}}}\left[ -\left( \frac{\overrightarrow{R}}{c}\frac{\partial \overrightarrow{a}}{\partial t'} \right)\frac{{{R}_{x}}{{R}_{x}}}{K}+3\left( \overrightarrow{v}\frac{\overrightarrow{a}}{c} \right)\frac{{{R}_{x}}{{R}_{x}}}{K}+\left( \frac{\overrightarrow{R}}{c}\frac{\partial \overrightarrow{a}}{\partial t'} \right)\frac{R}{c}\frac{{{R}_{x}}{{v}_{x}}}{K}-\frac{3R}{c}{{a}_{x}}{{v}_{x}}-3\frac{R}{c}\left( \overrightarrow{v}\frac{\overrightarrow{a}}{c} \right)\frac{{{R}_{x}}{{v}_{x}}}{K} \right]+ \\ 
 & +\left( 1+\frac{\overrightarrow{a}\cdot \overrightarrow{R}}{{{c}^{2}}}-\frac{\overrightarrow{v}\cdot \overrightarrow{v}}{{{c}^{2}}} \right)\frac{1}{{{K}^{3}}}\left[ 1+\frac{3R}{K}\left\{ \frac{{{v}_{x}}}{c}\frac{{{R}_{x}}}{K}\left( 1+\frac{\overrightarrow{R}\cdot \overrightarrow{a}-{{v}^{2}}}{{{c}^{2}}} \right)-\frac{{{v}_{x}}}{c}\frac{{{v}_{x}}}{c} \right\}-\frac{3}{K}\left\{ \frac{{{R}_{x}}{{R}_{x}}}{K}\left( 1+\frac{\overrightarrow{R}\cdot \overrightarrow{a}-{{v}^{2}}}{{{c}^{2}}} \right)-\frac{{{v}_{x}}{{R}_{x}}}{c} \right\}+\frac{R{{R}_{x}}}{K}\frac{{{a}_{x}}}{{{c}^{2}}} \right]+ \\ 
 & +\frac{1}{{{K}^{3}}{{c}^{2}}}\left( 2R\frac{{{a}_{x}}{{R}_{x}}}{K}\left( 1+\frac{\overrightarrow{R}\cdot \overrightarrow{a}-{{v}^{2}}}{{{c}^{2}}} \right)+\frac{R{{R}_{x}}}{c}\frac{\partial {{a}_{x}}}{\partial t'} \right) \\ 
\end{align}\]

\[\begin{align}
  & \frac{1}{e}\frac{\partial {{E}_{x}}}{\partial x}=\frac{1}{{{c}^{2}}}\frac{1}{{{K}^{3}}}\left[ -\left( \frac{\overrightarrow{R}}{c}\frac{\partial \overrightarrow{a}}{\partial t'} \right)\frac{{{R}_{x}}{{R}_{x}}}{K}+3\left( \overrightarrow{v}\frac{\overrightarrow{a}}{c} \right)\frac{{{R}_{x}}{{R}_{x}}}{K}+\left( \frac{\overrightarrow{R}}{c}\frac{\partial \overrightarrow{a}}{\partial t'} \right)\frac{R}{c}\frac{{{R}_{x}}{{v}_{x}}}{K}-\frac{3R}{c}{{a}_{x}}{{v}_{x}}-3\frac{R}{c}\left( \overrightarrow{v}\frac{\overrightarrow{a}}{c} \right)\frac{{{R}_{x}}{{v}_{x}}}{K} \right]+ \\ 
 & +\left( 1+\frac{\overrightarrow{a}\cdot \overrightarrow{R}}{{{c}^{2}}}-\frac{\overrightarrow{v}\cdot \overrightarrow{v}}{{{c}^{2}}} \right)\frac{1}{{{K}^{3}}}\left[ 1+\frac{2R}{{{c}^{2}}}\frac{{{a}_{x}}{{R}_{x}}}{K}+\frac{3R}{K}\left\{ \frac{{{v}_{x}}}{c}\frac{{{R}_{x}}}{K}\left( 1+\frac{\overrightarrow{R}\cdot \overrightarrow{a}-{{v}^{2}}}{{{c}^{2}}} \right)-\frac{{{v}_{x}}}{c}\frac{{{v}_{x}}}{c} \right\}-\frac{3}{K}\left\{ \frac{{{R}_{x}}{{R}_{x}}}{K}\left( 1+\frac{\overrightarrow{R}\cdot \overrightarrow{a}-{{v}^{2}}}{{{c}^{2}}} \right)-\frac{{{v}_{x}}{{R}_{x}}}{c} \right\}+\frac{R{{R}_{x}}}{K}\frac{{{a}_{x}}}{{{c}^{2}}} \right]+ \\ 
 & +\frac{1}{{{K}^{3}}{{c}^{2}}}\left( \frac{R{{R}_{x}}}{c}\frac{\partial {{a}_{x}}}{\partial t'} \right) \\ 
\end{align}\]

\[\begin{align}
  & \frac{1}{e}\frac{\partial {{E}_{x}}}{\partial x}=\frac{1}{{{c}^{2}}}\frac{1}{{{K}^{3}}}\left[ -\left( \frac{\overrightarrow{R}}{c}\frac{\partial \overrightarrow{a}}{\partial t'} \right)\frac{{{R}_{x}}{{R}_{x}}}{K}+3\left( \overrightarrow{v}\frac{\overrightarrow{a}}{c} \right)\frac{{{R}_{x}}{{R}_{x}}}{K}+\left( \frac{\overrightarrow{R}}{c}\frac{\partial \overrightarrow{a}}{\partial t'} \right)\frac{R}{c}\frac{{{R}_{x}}{{v}_{x}}}{K}-\frac{3R}{c}{{a}_{x}}{{v}_{x}}-3\frac{R}{c}\left( \overrightarrow{v}\frac{\overrightarrow{a}}{c} \right)\frac{{{R}_{x}}{{v}_{x}}}{K}+\frac{R{{R}_{x}}}{c}\frac{\partial {{a}_{x}}}{\partial t'} \right]+ \\ 
 & +\left( 1+\frac{\overrightarrow{a}\cdot \overrightarrow{R}}{{{c}^{2}}}-\frac{\overrightarrow{v}\cdot \overrightarrow{v}}{{{c}^{2}}} \right)\frac{1}{{{K}^{3}}}\left[ 1+\frac{2R}{{{c}^{2}}}\frac{{{a}_{x}}{{R}_{x}}}{K}+\frac{3R}{K}\left\{ \frac{{{v}_{x}}}{c}\frac{{{R}_{x}}}{K}\left( 1+\frac{\overrightarrow{R}\cdot \overrightarrow{a}-{{v}^{2}}}{{{c}^{2}}} \right)-\frac{{{v}_{x}}}{c}\frac{{{v}_{x}}}{c} \right\}-\frac{3}{K}\left\{ \frac{{{R}_{x}}{{R}_{x}}}{K}\left( 1+\frac{\overrightarrow{R}\cdot \overrightarrow{a}-{{v}^{2}}}{{{c}^{2}}} \right)-\frac{{{v}_{x}}{{R}_{x}}}{c} \right\}+\frac{R{{R}_{x}}}{K}\frac{{{a}_{x}}}{{{c}^{2}}} \right] \\ 
\end{align}\] Выражения производных \(y\) и \(z\) компонент
электрического поля по \(y\)и \(z\) координатам точки наблюдения легко
получить аналогичной заменой индексов в данном выражении Скалярно
складывая производные\(x\) , \(y\) и \(z\)компонент электрического поля
по соответственно \(x\) , \(y\) и \(z\)координатам точки наблюдения
получаем выражение для дивергенции электрического поля \[\begin{align}
  & \frac{1}{e}\left( \frac{\partial {{E}_{x}}}{\partial x}+\frac{\partial {{E}_{y}}}{\partial y}+\frac{\partial {{E}_{z}}}{\partial z} \right)=\frac{1}{{{c}^{2}}}\frac{1}{{{K}^{3}}}\left[ -\left( \frac{\overrightarrow{R}}{c}\frac{\partial \overrightarrow{a}}{\partial t'} \right)\frac{{{R}^{2}}}{K}+3\left( \overrightarrow{v}\frac{\overrightarrow{a}}{c} \right)\frac{{{R}^{2}}}{K}+\left( \frac{\overrightarrow{R}}{c}\frac{\partial \overrightarrow{a}}{\partial t'} \right)\frac{R}{c}\frac{\overrightarrow{R}\overrightarrow{v}}{K}-\frac{3R}{c}\overrightarrow{a}\overrightarrow{v}-3\frac{R}{c}\left( \overrightarrow{v}\frac{\overrightarrow{a}}{c} \right)\frac{\overrightarrow{R}\overrightarrow{v}}{K}+\frac{R\overrightarrow{R}}{c}\frac{\partial \overrightarrow{a}}{\partial t'} \right]+ \\ 
 & +\left( 1+\frac{\overrightarrow{a}\cdot \overrightarrow{R}}{{{c}^{2}}}-\frac{\overrightarrow{v}\cdot \overrightarrow{v}}{{{c}^{2}}} \right)\frac{1}{{{K}^{3}}}\left[ 3+\frac{2R}{{{c}^{2}}}\frac{\overrightarrow{a}\overrightarrow{R}}{K}+\frac{3R}{K}\left\{ \frac{\overrightarrow{v}}{c}\frac{\overrightarrow{R}}{K}\left( 1+\frac{\overrightarrow{R}\cdot \overrightarrow{a}-{{v}^{2}}}{{{c}^{2}}} \right)-\frac{\overrightarrow{v}}{c}\frac{\overrightarrow{v}}{c} \right\}-\frac{3}{K}\left\{ \frac{{{R}^{2}}}{K}\left( 1+\frac{\overrightarrow{R}\cdot \overrightarrow{a}-{{v}^{2}}}{{{c}^{2}}} \right)-\frac{\overrightarrow{v}\overrightarrow{R}}{c} \right\}+\frac{R\overrightarrow{R}}{K}\frac{\overrightarrow{a}}{{{c}^{2}}} \right] \\ 
 &  \\ 
\end{align}\]

\[\begin{align}
  & \frac{1}{e}\left( \frac{\partial {{E}_{x}}}{\partial x}+\frac{\partial {{E}_{y}}}{\partial y}+\frac{\partial {{E}_{z}}}{\partial z} \right)=\frac{1}{{{c}^{2}}}\frac{1}{{{K}^{3}}}\left[ -\left( \frac{\overrightarrow{R}}{c}\frac{\partial \overrightarrow{a}}{\partial t'} \right)\frac{{{R}^{2}}}{K}+3\left( \overrightarrow{v}\frac{\overrightarrow{a}}{c} \right)\frac{{{R}^{2}}}{K}+\left( \frac{\overrightarrow{R}}{c}\frac{\partial \overrightarrow{a}}{\partial t'} \right)\frac{R}{c}\frac{\overrightarrow{R}\overrightarrow{v}}{K}-\frac{3R}{c}\overrightarrow{a}\overrightarrow{v}-3\frac{R}{c}\left( \overrightarrow{v}\frac{\overrightarrow{a}}{c} \right)\frac{\overrightarrow{R}\overrightarrow{v}}{K}+\frac{R\overrightarrow{R}}{c}\frac{\partial \overrightarrow{a}}{\partial t'} \right]+ \\ 
 & +\left( 1+\frac{\overrightarrow{a}\cdot \overrightarrow{R}}{{{c}^{2}}}-\frac{\overrightarrow{v}\cdot \overrightarrow{v}}{{{c}^{2}}} \right)\frac{1}{{{K}^{3}}}\left[ 3+\frac{2R}{{{c}^{2}}}\frac{\overrightarrow{a}\overrightarrow{R}}{K}+\frac{R\overrightarrow{R}}{K}\frac{\overrightarrow{a}}{{{c}^{2}}}+\frac{3R}{K}\left\{ \frac{\overrightarrow{v}}{c}\frac{\overrightarrow{R}}{K}\left( 1+\frac{\overrightarrow{R}\cdot \overrightarrow{a}-{{v}^{2}}}{{{c}^{2}}} \right)-\frac{\overrightarrow{v}}{c}\frac{\overrightarrow{v}}{c} \right\}-\frac{3}{K}\left\{ \frac{{{R}^{2}}}{K}\left( 1+\frac{\overrightarrow{R}\cdot \overrightarrow{a}-{{v}^{2}}}{{{c}^{2}}} \right)-\frac{\overrightarrow{v}\overrightarrow{R}}{c} \right\} \right] \\ 
 &  \\ 
\end{align}\]

\[\begin{align}
  & \frac{1}{e}\left( \frac{\partial {{E}_{x}}}{\partial x}+\frac{\partial {{E}_{y}}}{\partial y}+\frac{\partial {{E}_{z}}}{\partial z} \right)=\frac{1}{{{c}^{2}}}\frac{1}{{{K}^{3}}}\left[ \left( \frac{\overrightarrow{R}}{c}\frac{\partial \overrightarrow{a}}{\partial t'} \right)\frac{R}{c}\frac{\overrightarrow{R}\overrightarrow{v}}{K}-\left( \frac{\overrightarrow{R}}{c}\frac{\partial \overrightarrow{a}}{\partial t'} \right)\frac{{{R}^{2}}}{K}+\frac{R\overrightarrow{R}}{c}\frac{\partial \overrightarrow{a}}{\partial t'}+3\left( \overrightarrow{v}\frac{\overrightarrow{a}}{c} \right)\frac{{{R}^{2}}}{K}-\frac{3R}{c}\overrightarrow{a}\overrightarrow{v}-3\frac{R}{c}\left( \overrightarrow{v}\frac{\overrightarrow{a}}{c} \right)\frac{\overrightarrow{R}\overrightarrow{v}}{K} \right]+ \\ 
 & +\left( 1+\frac{\overrightarrow{a}\cdot \overrightarrow{R}}{{{c}^{2}}}-\frac{\overrightarrow{v}\cdot \overrightarrow{v}}{{{c}^{2}}} \right)\frac{1}{{{K}^{3}}}\left[ 3+\frac{2R}{{{c}^{2}}}\frac{\overrightarrow{a}\overrightarrow{R}}{K}+\frac{R\overrightarrow{R}}{K}\frac{\overrightarrow{a}}{{{c}^{2}}}+\frac{3R}{K}\frac{\overrightarrow{v}}{c}\frac{\overrightarrow{R}}{K}\left( 1+\frac{\overrightarrow{R}\cdot \overrightarrow{a}-{{v}^{2}}}{{{c}^{2}}} \right)-\frac{3R}{K}\frac{\overrightarrow{v}}{c}\frac{\overrightarrow{v}}{c}-\frac{3}{K}\frac{{{R}^{2}}}{K}\left( 1+\frac{\overrightarrow{R}\cdot \overrightarrow{a}-{{v}^{2}}}{{{c}^{2}}} \right)+\frac{3}{K}\frac{\overrightarrow{v}\overrightarrow{R}}{c} \right] \\ 
 &  \\ 
\end{align}\]

\[\begin{align}
  & \frac{1}{e}\left( \frac{\partial {{E}_{x}}}{\partial x}+\frac{\partial {{E}_{y}}}{\partial y}+\frac{\partial {{E}_{z}}}{\partial z} \right)=\frac{1}{{{c}^{2}}}\frac{1}{{{K}^{3}}}\left[ \left( \frac{\overrightarrow{R}}{c}\frac{\partial \overrightarrow{a}}{\partial t'} \right)\left\{ \frac{R}{c}\frac{\overrightarrow{R}\overrightarrow{v}}{K}-\frac{{{R}^{2}}}{K}+R \right\}+3\left( \frac{\overrightarrow{a}\overrightarrow{v}}{c} \right)\left\{ \frac{{{R}^{2}}}{K}-R-\frac{R}{c}\frac{\overrightarrow{R}\overrightarrow{v}}{K} \right\} \right]+ \\ 
 & +\left( 1+\frac{\overrightarrow{a}\cdot \overrightarrow{R}}{{{c}^{2}}}-\frac{\overrightarrow{v}\cdot \overrightarrow{v}}{{{c}^{2}}} \right)\frac{1}{{{K}^{3}}}\left[ 3+3\frac{R}{{{c}^{2}}}\frac{\overrightarrow{a}\overrightarrow{R}}{K}+\frac{3R}{K}\frac{\overrightarrow{v}}{c}\frac{\overrightarrow{R}}{K}\left( 1+\frac{\overrightarrow{R}\cdot \overrightarrow{a}-{{v}^{2}}}{{{c}^{2}}} \right)-\frac{3R}{K}\frac{\overrightarrow{v}}{c}\frac{\overrightarrow{v}}{c}-\frac{3}{K}\frac{{{R}^{2}}}{K}\left( 1+\frac{\overrightarrow{R}\cdot \overrightarrow{a}-{{v}^{2}}}{{{c}^{2}}} \right)+\frac{3}{K}\frac{\overrightarrow{v}\overrightarrow{R}}{c} \right] \\ 
 &  \\ 
\end{align}\]
\[\left\{ \frac{R}{c}\frac{\overrightarrow{R}\overrightarrow{v}}{K}-\frac{{{R}^{2}}}{K}+R \right\}=\left\{ \frac{R\left( \frac{\overrightarrow{v}\cdot \overrightarrow{R}}{c} \right)}{R-\frac{\overrightarrow{v}\cdot \overrightarrow{R}}{c}}-\frac{{{R}^{2}}}{R-\frac{\overrightarrow{v}\cdot \overrightarrow{R}}{c}}+\frac{R\left( R-\frac{\overrightarrow{v}\cdot \overrightarrow{R}}{c} \right)}{R-\frac{\overrightarrow{v}\cdot \overrightarrow{R}}{c}} \right\}=\frac{R\left( \frac{\overrightarrow{v}\cdot \overrightarrow{R}}{c} \right)-{{R}^{2}}+R\left( R-\frac{\overrightarrow{v}\cdot \overrightarrow{R}}{c} \right)}{R-\frac{\overrightarrow{v}\cdot \overrightarrow{R}}{c}}=0\]
\[\begin{align}
  & \frac{1}{e}\left( \frac{\partial {{E}_{x}}}{\partial x}+\frac{\partial {{E}_{y}}}{\partial y}+\frac{\partial {{E}_{z}}}{\partial z} \right)=\left( 1+\frac{\overrightarrow{a}\cdot \overrightarrow{R}}{{{c}^{2}}}-\frac{\overrightarrow{v}\cdot \overrightarrow{v}}{{{c}^{2}}} \right)\frac{1}{{{K}^{3}}} \\ 
 & \left[ 3+3\frac{R}{{{c}^{2}}}\frac{\overrightarrow{a}\overrightarrow{R}}{K}-\frac{3R}{K}\frac{\overrightarrow{v}}{c}\frac{\overrightarrow{v}}{c}+\frac{3R}{K}\frac{\overrightarrow{v}}{c}\frac{\overrightarrow{R}}{K}\left( 1+\frac{\overrightarrow{R}\cdot \overrightarrow{a}-{{v}^{2}}}{{{c}^{2}}} \right)-\frac{3}{K}\frac{{{R}^{2}}}{K}\left( 1+\frac{\overrightarrow{R}\cdot \overrightarrow{a}-{{v}^{2}}}{{{c}^{2}}} \right)+\frac{3}{K}\frac{\overrightarrow{v}\overrightarrow{R}}{c} \right] \\ 
\end{align}\]

\[\begin{align}
  & \frac{1}{e}\left( \frac{\partial {{E}_{x}}}{\partial x}+\frac{\partial {{E}_{y}}}{\partial y}+\frac{\partial {{E}_{z}}}{\partial z} \right)=\frac{1}{{{K}^{3}}}\left( 1+\frac{\overrightarrow{a}\cdot \overrightarrow{R}}{{{c}^{2}}}-\frac{\overrightarrow{v}\cdot \overrightarrow{v}}{{{c}^{2}}} \right) \\ 
 & \left[ 3+3\frac{R}{K}\frac{\overrightarrow{a}\cdot \overrightarrow{R}}{{{c}^{2}}}-\frac{3R}{K}\frac{{{v}^{2}}}{{{c}^{2}}}+3\left\{ \frac{R}{{{K}^{2}}}\left( \frac{\overrightarrow{v}\cdot \overrightarrow{R}}{c} \right)-\frac{{{R}^{2}}}{{{K}^{2}}} \right\}\left( 1+\frac{\overrightarrow{R}\cdot \overrightarrow{a}-{{v}^{2}}}{{{c}^{2}}} \right)+\frac{3}{K}\left( \frac{\overrightarrow{v}\cdot \overrightarrow{R}}{c} \right) \right] \\ 
\end{align}\]
\[\left\{ \frac{R}{{{K}^{2}}}\left( \frac{\overrightarrow{v}\cdot \overrightarrow{R}}{c} \right)-\frac{{{R}^{2}}}{{{K}^{2}}} \right\}=\frac{R}{K}\left\{ \frac{1}{K}\left( \frac{\overrightarrow{v}\cdot \overrightarrow{R}}{c} \right)-\frac{R}{K} \right\}=\frac{R}{K}\left\{ \frac{\frac{\overrightarrow{v}\cdot \overrightarrow{R}}{c}}{R-\frac{\overrightarrow{v}\cdot \overrightarrow{R}}{c}}-\frac{R}{R-\frac{\overrightarrow{v}\cdot \overrightarrow{R}}{c}} \right\}=\frac{R}{K}\left\{ \frac{\frac{\overrightarrow{v}\cdot \overrightarrow{R}}{c}-R}{R-\frac{\overrightarrow{v}\cdot \overrightarrow{R}}{c}} \right\}=-\frac{R}{K}\]
Итак, дивергенция электрического поля равна
\[\frac{1}{e}\left( \frac{\partial {{E}_{x}}}{\partial x}+\frac{\partial {{E}_{y}}}{\partial y}+\frac{\partial {{E}_{z}}}{\partial z} \right)=\frac{1}{{{K}^{3}}}\left( 1+\frac{\overrightarrow{a}\cdot \overrightarrow{R}}{{{c}^{2}}}-\frac{\overrightarrow{v}\cdot \overrightarrow{v}}{{{c}^{2}}} \right)\left[ 3+3\frac{R}{K}\frac{\overrightarrow{a}\cdot \overrightarrow{R}}{{{c}^{2}}}-\frac{3R}{K}\frac{{{v}^{2}}}{{{c}^{2}}}-3\frac{R}{K}\left( 1+\frac{\overrightarrow{R}\cdot \overrightarrow{a}-{{v}^{2}}}{{{c}^{2}}} \right)+\frac{3}{K}\left( \frac{\overrightarrow{v}\cdot \overrightarrow{R}}{c} \right) \right]\]
При этом дивергенция первого слагаемого электрического поля
\[\begin{align}
  & \frac{1}{e}{{\left( \frac{\partial {{E}_{x}}}{\partial x}+\frac{\partial {{E}_{y}}}{\partial y}+\frac{\partial {{E}_{z}}}{\partial z} \right)}_{1}}=-\frac{{{R}^{2}}}{{{K}^{4}}{{c}^{2}}}\left( \frac{\overrightarrow{R}}{c}\frac{\partial \overrightarrow{a}}{\partial t'} \right)+\frac{{{R}^{2}}}{{{K}^{4}}{{c}^{2}}}\left( \frac{3\ \overrightarrow{v}\cdot \overrightarrow{a}}{c} \right)-\frac{2}{{{K}^{3}}}+ \\ 
 & +\left( 1+\frac{\overrightarrow{a}\cdot \overrightarrow{R}}{{{c}^{2}}}-\frac{\overrightarrow{v}\cdot \overrightarrow{v}}{{{c}^{2}}} \right)\cdot \left\{ \frac{5}{{{K}^{3}}}-\frac{3{{R}^{2}}}{{{K}^{5}}}\left( 1+\frac{\overrightarrow{R}\cdot \overrightarrow{a}-{{v}^{2}}}{{{c}^{2}}} \right)+6\frac{\overrightarrow{v}\cdot \overrightarrow{R}}{c{{K}^{4}}} \right\} \\ 
\end{align}\] И дивергенция второго слагаемого \[\begin{align}
  & \frac{1}{e}{{\left( \frac{\partial {{E}_{x}}}{\partial x}+\frac{\partial {{E}_{y}}}{\partial y}+\frac{\partial {{E}_{z}}}{\partial z} \right)}_{2}}=\frac{R}{{{K}^{3}}}\frac{1}{{{c}^{2}}}\left\{ \left( \frac{\overrightarrow{R}}{c}\frac{\partial \overrightarrow{a}}{\partial t'} \right)\left( \frac{\overrightarrow{R}}{K}\frac{\overrightarrow{v}}{c}+1 \right)-3\left( \overrightarrow{v}\frac{\overrightarrow{a}}{c} \right)\left( \frac{\overrightarrow{R}}{K}\frac{\overrightarrow{v}}{c}+1 \right) \right\}- \\ 
 & -\left( 1+\frac{\overrightarrow{R}\cdot \overrightarrow{a}-{{v}^{2}}}{{{c}^{2}}} \right)\left( 3\frac{\overrightarrow{R}}{{{K}^{4}}}\frac{\overrightarrow{v}}{c}-\frac{3R}{{{K}^{4}}}\left\{ \frac{\overrightarrow{R}}{K}\frac{\overrightarrow{v}}{c}\left( 1+\frac{\overrightarrow{R}\cdot \overrightarrow{a}-{{v}^{2}}}{{{c}^{2}}} \right)-\frac{\overrightarrow{v}}{c}\frac{\overrightarrow{v}}{c} \right\}-\frac{3R\overrightarrow{R}}{{{K}^{4}}{{c}^{2}}}\overrightarrow{a} \right)-2\frac{\overrightarrow{R}}{{{K}^{3}}}\frac{\overrightarrow{a}}{{{c}^{2}}}+\frac{2}{{{K}^{3}}}\frac{\overrightarrow{v}}{c}\frac{\overrightarrow{v}}{c} \\ 
\end{align}\] Для проверки сумма дивергенций обоих слагаемых
электрического поля \[\begin{align}
  & \frac{1}{e}\left( \frac{\partial {{E}_{x}}}{\partial x}+\frac{\partial {{E}_{y}}}{\partial y}+\frac{\partial {{E}_{z}}}{\partial z} \right)=-\frac{R}{{{K}^{3}}{{c}^{2}}}\frac{R}{K}\left( \frac{\overrightarrow{R}}{c}\frac{\partial \overrightarrow{a}}{\partial t'} \right)+\frac{R}{{{K}^{3}}{{c}^{2}}}\frac{R}{K}\left( \frac{3\ \overrightarrow{v}\cdot \overrightarrow{a}}{c} \right)-\frac{2}{{{K}^{3}}}+ \\ 
 & +\frac{R}{{{K}^{3}}}\frac{1}{{{c}^{2}}}\left\{ \left( \frac{\overrightarrow{R}}{c}\frac{\partial \overrightarrow{a}}{\partial t'} \right)\left( \frac{\overrightarrow{R}}{K}\frac{\overrightarrow{v}}{c}+1 \right)-3\left( \overrightarrow{v}\frac{\overrightarrow{a}}{c} \right)\left( \frac{\overrightarrow{R}}{K}\frac{\overrightarrow{v}}{c}+1 \right) \right\}+ \\ 
 & +\left( 1+\frac{\overrightarrow{a}\cdot \overrightarrow{R}}{{{c}^{2}}}-\frac{\overrightarrow{v}\cdot \overrightarrow{v}}{{{c}^{2}}} \right)\cdot \left\{ \frac{5}{{{K}^{3}}}-\frac{3{{R}^{2}}}{{{K}^{5}}}\left( 1+\frac{\overrightarrow{R}\cdot \overrightarrow{a}-{{v}^{2}}}{{{c}^{2}}} \right)+6\frac{\overrightarrow{v}\cdot \overrightarrow{R}}{c{{K}^{4}}} \right\}- \\ 
 & -\left( 1+\frac{\overrightarrow{R}\cdot \overrightarrow{a}-{{v}^{2}}}{{{c}^{2}}} \right)\left( 3\frac{\overrightarrow{R}}{{{K}^{4}}}\frac{\overrightarrow{v}}{c}-\frac{3R}{{{K}^{4}}}\left\{ \frac{\overrightarrow{R}}{K}\frac{\overrightarrow{v}}{c}\left( 1+\frac{\overrightarrow{R}\cdot \overrightarrow{a}-{{v}^{2}}}{{{c}^{2}}} \right)-\frac{\overrightarrow{v}}{c}\frac{\overrightarrow{v}}{c} \right\}-\frac{3R\overrightarrow{R}}{{{K}^{4}}{{c}^{2}}}\overrightarrow{a} \right)-2\frac{\overrightarrow{R}}{{{K}^{3}}}\frac{\overrightarrow{a}}{{{c}^{2}}}+\frac{2}{{{K}^{3}}}\frac{\overrightarrow{v}}{c}\frac{\overrightarrow{v}}{c} \\ 
\end{align}\]

\[\begin{align}
  & \frac{1}{e}\left( \frac{\partial {{E}_{x}}}{\partial x}+\frac{\partial {{E}_{y}}}{\partial y}+\frac{\partial {{E}_{z}}}{\partial z} \right)=-\frac{2}{{{K}^{3}}}-2\frac{\overrightarrow{R}}{{{K}^{3}}}\frac{\overrightarrow{a}}{{{c}^{2}}}+\frac{2}{{{K}^{3}}}\frac{\overrightarrow{v}}{c}\frac{\overrightarrow{v}}{c}+ \\ 
 & +\frac{R}{{{K}^{3}}}\frac{1}{{{c}^{2}}}\left\{ \left( \frac{\overrightarrow{R}}{c}\frac{\partial \overrightarrow{a}}{\partial t'} \right)\left( \frac{\overrightarrow{R}}{K}\frac{\overrightarrow{v}}{c}-\frac{R}{K}+1 \right)-3\left( \overrightarrow{v}\frac{\overrightarrow{a}}{c} \right)\left( \frac{\overrightarrow{R}}{K}\frac{\overrightarrow{v}}{c}-\frac{R}{K}+1 \right) \right\}+ \\ 
 & +\left( 1+\frac{\overrightarrow{a}\cdot \overrightarrow{R}}{{{c}^{2}}}-\frac{\overrightarrow{v}\cdot \overrightarrow{v}}{{{c}^{2}}} \right)\cdot \left\{ \frac{5}{{{K}^{3}}}-\frac{3{{R}^{2}}}{{{K}^{5}}}\left( 1+\frac{\overrightarrow{R}\cdot \overrightarrow{a}-{{v}^{2}}}{{{c}^{2}}} \right)+6\frac{\overrightarrow{v}\cdot \overrightarrow{R}}{c{{K}^{4}}}-3\frac{\overrightarrow{R}}{{{K}^{4}}}\frac{\overrightarrow{v}}{c}+\frac{3R}{{{K}^{4}}}\left\{ \frac{\overrightarrow{R}}{K}\frac{\overrightarrow{v}}{c}\left( 1+\frac{\overrightarrow{R}\cdot \overrightarrow{a}-{{v}^{2}}}{{{c}^{2}}} \right)-\frac{\overrightarrow{v}}{c}\frac{\overrightarrow{v}}{c} \right\}+\frac{3R\overrightarrow{R}}{{{K}^{4}}{{c}^{2}}}\overrightarrow{a} \right\} \\ 
 &  \\ 
\end{align}\]
\[\left\{ \frac{1}{c}\frac{\overrightarrow{R}\overrightarrow{v}}{K}-\frac{R}{K}+1 \right\}=0\]
\[\begin{align}
  & \frac{1}{e}\left( \frac{\partial {{E}_{x}}}{\partial x}+\frac{\partial {{E}_{y}}}{\partial y}+\frac{\partial {{E}_{z}}}{\partial z} \right)=-\frac{2}{{{K}^{3}}}\left( 1+\frac{\overrightarrow{a}\cdot \overrightarrow{R}}{{{c}^{2}}}-\frac{\overrightarrow{v}\cdot \overrightarrow{v}}{{{c}^{2}}} \right)+ \\ 
 & +\left( 1+\frac{\overrightarrow{a}\cdot \overrightarrow{R}}{{{c}^{2}}}-\frac{\overrightarrow{v}\cdot \overrightarrow{v}}{{{c}^{2}}} \right)\frac{1}{{{K}^{3}}}\left\{ 5-\frac{3{{R}^{2}}}{{{K}^{2}}}\left( 1+\frac{\overrightarrow{R}\cdot \overrightarrow{a}-{{v}^{2}}}{{{c}^{2}}} \right)+3\frac{\overrightarrow{v}\cdot \overrightarrow{R}}{cK}+\frac{3R}{K}\left\{ \frac{\overrightarrow{R}}{K}\frac{\overrightarrow{v}}{c}\left( 1+\frac{\overrightarrow{R}\cdot \overrightarrow{a}-{{v}^{2}}}{{{c}^{2}}} \right)-\frac{\overrightarrow{v}}{c}\frac{\overrightarrow{v}}{c} \right\}+\frac{3R\overrightarrow{R}}{K{{c}^{2}}}\overrightarrow{a} \right\} \\ 
\end{align}\]

\[\frac{1}{e}\left( \frac{\partial {{E}_{x}}}{\partial x}+\frac{\partial {{E}_{y}}}{\partial y}+\frac{\partial {{E}_{z}}}{\partial z} \right)=\left( 1+\frac{\overrightarrow{a}\cdot \overrightarrow{R}}{{{c}^{2}}}-\frac{\overrightarrow{v}\cdot \overrightarrow{v}}{{{c}^{2}}} \right)\frac{1}{{{K}^{3}}}\left\{ 3-\frac{3{{R}^{2}}}{{{K}^{2}}}\left( 1+\frac{\overrightarrow{R}\cdot \overrightarrow{a}-{{v}^{2}}}{{{c}^{2}}} \right)+3\frac{\overrightarrow{v}\cdot \overrightarrow{R}}{cK}+\frac{3R}{K}\left\{ \frac{\overrightarrow{R}}{K}\frac{\overrightarrow{v}}{c}\left( 1+\frac{\overrightarrow{R}\cdot \overrightarrow{a}-{{v}^{2}}}{{{c}^{2}}} \right)-\frac{\overrightarrow{v}}{c}\frac{\overrightarrow{v}}{c} \right\}+\frac{3R\overrightarrow{R}}{K{{c}^{2}}}\overrightarrow{a} \right\}\]

\[\frac{1}{e}\left( \frac{\partial {{E}_{x}}}{\partial x}+\frac{\partial {{E}_{y}}}{\partial y}+\frac{\partial {{E}_{z}}}{\partial z} \right)=\left( 1+\frac{\overrightarrow{a}\cdot \overrightarrow{R}}{{{c}^{2}}}-\frac{\overrightarrow{v}\cdot \overrightarrow{v}}{{{c}^{2}}} \right)\frac{1}{{{K}^{3}}}\left\{ 3-\frac{3{{R}^{2}}}{{{K}^{2}}}\left( 1+\frac{\overrightarrow{R}\cdot \overrightarrow{a}-{{v}^{2}}}{{{c}^{2}}} \right)+3\frac{\overrightarrow{v}\cdot \overrightarrow{R}}{cK}+\frac{3R}{K}\frac{\overrightarrow{R}}{K}\frac{\overrightarrow{v}}{c}\left( 1+\frac{\overrightarrow{R}\cdot \overrightarrow{a}-{{v}^{2}}}{{{c}^{2}}} \right)-\frac{3R}{K}\frac{\overrightarrow{v}}{c}\frac{\overrightarrow{v}}{c}+3\frac{R}{K}\frac{\overrightarrow{R}\cdot \overrightarrow{a}}{{{c}^{2}}} \right\}\]
\[\frac{1}{e}\left( \frac{\partial {{E}_{x}}}{\partial x}+\frac{\partial {{E}_{y}}}{\partial y}+\frac{\partial {{E}_{z}}}{\partial z} \right)=\left( 1+\frac{\overrightarrow{a}\cdot \overrightarrow{R}}{{{c}^{2}}}-\frac{\overrightarrow{v}\cdot \overrightarrow{v}}{{{c}^{2}}} \right)\frac{1}{{{K}^{3}}}\left\{ 3-\frac{3R}{K}\frac{R}{K}\left( 1+\frac{\overrightarrow{R}\cdot \overrightarrow{a}-{{v}^{2}}}{{{c}^{2}}} \right)+\frac{3R}{K}\frac{\overrightarrow{R}}{K}\frac{\overrightarrow{v}}{c}\left( 1+\frac{\overrightarrow{R}\cdot \overrightarrow{a}-{{v}^{2}}}{{{c}^{2}}} \right)+3\frac{\overrightarrow{v}\cdot \overrightarrow{R}}{cK}-\frac{3R}{K}\frac{\overrightarrow{v}}{c}\frac{\overrightarrow{v}}{c}+3\frac{R}{K}\frac{\overrightarrow{R}\cdot \overrightarrow{a}}{{{c}^{2}}} \right\}\]
\[\frac{1}{e}\left( \frac{\partial {{E}_{x}}}{\partial x}+\frac{\partial {{E}_{y}}}{\partial y}+\frac{\partial {{E}_{z}}}{\partial z} \right)=\left( 1+\frac{\overrightarrow{a}\cdot \overrightarrow{R}}{{{c}^{2}}}-\frac{\overrightarrow{v}\cdot \overrightarrow{v}}{{{c}^{2}}} \right)\frac{1}{{{K}^{3}}}\left\{ 3+\frac{3R}{K}\left( \frac{\overrightarrow{R}}{K}\frac{\overrightarrow{v}}{c}-\frac{R}{K} \right)\left( 1+\frac{\overrightarrow{R}\cdot \overrightarrow{a}-{{v}^{2}}}{{{c}^{2}}} \right)+3\frac{\overrightarrow{v}\cdot \overrightarrow{R}}{cK}-\frac{3R}{K}\frac{\overrightarrow{v}}{c}\frac{\overrightarrow{v}}{c}+3\frac{R}{K}\frac{\overrightarrow{R}\cdot \overrightarrow{a}}{{{c}^{2}}} \right\}\]
\[\left\{ \frac{1}{K}\left( \frac{\overrightarrow{v}\cdot \overrightarrow{R}}{c} \right)-\frac{R}{K} \right\}=-1\]
\[\frac{1}{e}\left( \frac{\partial {{E}_{x}}}{\partial x}+\frac{\partial {{E}_{y}}}{\partial y}+\frac{\partial {{E}_{z}}}{\partial z} \right)=\left( 1+\frac{\overrightarrow{a}\cdot \overrightarrow{R}}{{{c}^{2}}}-\frac{\overrightarrow{v}\cdot \overrightarrow{v}}{{{c}^{2}}} \right)\frac{1}{{{K}^{3}}}\left\{ 3-\frac{3R}{K}\left( 1+\frac{\overrightarrow{R}\cdot \overrightarrow{a}-{{v}^{2}}}{{{c}^{2}}} \right)+3\frac{\overrightarrow{v}\cdot \overrightarrow{R}}{cK}-\frac{3R}{K}\frac{\overrightarrow{v}}{c}\frac{\overrightarrow{v}}{c}+3\frac{R}{K}\frac{\overrightarrow{R}\cdot \overrightarrow{a}}{{{c}^{2}}} \right\}\]

\[\frac{1}{e}\left( \frac{\partial {{E}_{x}}}{\partial x}+\frac{\partial {{E}_{y}}}{\partial y}+\frac{\partial {{E}_{z}}}{\partial z} \right)=\frac{1}{{{K}^{3}}}\left( 1+\frac{\overrightarrow{a}\cdot \overrightarrow{R}}{{{c}^{2}}}-\frac{\overrightarrow{v}\cdot \overrightarrow{v}}{{{c}^{2}}} \right)\left[ 3+3\frac{R}{K}\frac{\overrightarrow{a}\cdot \overrightarrow{R}}{{{c}^{2}}}-\frac{3R}{K}\frac{{{v}^{2}}}{{{c}^{2}}}-3\frac{R}{K}\left( 1+\frac{\overrightarrow{R}\cdot \overrightarrow{a}-{{v}^{2}}}{{{c}^{2}}} \right)+\frac{3}{K}\left( \frac{\overrightarrow{v}\cdot \overrightarrow{R}}{c} \right) \right]\]

Сокращаем

\[\frac{1}{e}\left( \frac{\partial {{E}_{x}}}{\partial x}+\frac{\partial {{E}_{y}}}{\partial y}+\frac{\partial {{E}_{z}}}{\partial z} \right)=\frac{3}{{{K}^{3}}}\left( 1+\frac{\overrightarrow{a}\cdot \overrightarrow{R}}{{{c}^{2}}}-\frac{\overrightarrow{v}\cdot \overrightarrow{v}}{{{c}^{2}}} \right)\left[ 1+\frac{R}{K}\frac{\overrightarrow{a}\cdot \overrightarrow{R}}{{{c}^{2}}}-\frac{R}{K}\frac{{{v}^{2}}}{{{c}^{2}}}-\frac{R}{K}\left( 1+\frac{\overrightarrow{R}\cdot \overrightarrow{a}-{{v}^{2}}}{{{c}^{2}}} \right)+\frac{1}{K}\left( \frac{\overrightarrow{v}\cdot \overrightarrow{R}}{c} \right) \right]\]

\[\frac{1}{e}\left( \frac{\partial {{E}_{x}}}{\partial x}+\frac{\partial {{E}_{y}}}{\partial y}+\frac{\partial {{E}_{z}}}{\partial z} \right)=\frac{3}{{{K}^{3}}}\left( 1+\frac{\overrightarrow{a}\cdot \overrightarrow{R}}{{{c}^{2}}}-\frac{\overrightarrow{v}\cdot \overrightarrow{v}}{{{c}^{2}}} \right)\left[ 1+\frac{R}{K}\frac{\overrightarrow{a}\cdot \overrightarrow{R}}{{{c}^{2}}}-\frac{R}{K}\frac{{{v}^{2}}}{{{c}^{2}}}-\frac{R}{K}-\frac{R}{K}\frac{\overrightarrow{R}\cdot \overrightarrow{a}-{{v}^{2}}}{{{c}^{2}}}+\frac{1}{K}\left( \frac{\overrightarrow{v}\cdot \overrightarrow{R}}{c} \right) \right]\]

\[\frac{1}{e}\left( \frac{\partial {{E}_{x}}}{\partial x}+\frac{\partial {{E}_{y}}}{\partial y}+\frac{\partial {{E}_{z}}}{\partial z} \right)=\frac{3}{{{K}^{3}}}\left( 1+\frac{\overrightarrow{a}\cdot \overrightarrow{R}}{{{c}^{2}}}-\frac{\overrightarrow{v}\cdot \overrightarrow{v}}{{{c}^{2}}} \right)\left[ 1+\frac{R}{K}\frac{\overrightarrow{a}\cdot \overrightarrow{R}}{{{c}^{2}}}-\frac{R}{K}\frac{{{v}^{2}}}{{{c}^{2}}}-\frac{R}{K}\frac{\overrightarrow{R}\cdot \overrightarrow{a}-{{v}^{2}}}{{{c}^{2}}}+\frac{1}{K}\left( \frac{\overrightarrow{v}\cdot \overrightarrow{R}}{c} \right)-\frac{R}{K} \right]\]
\[\left\{ \frac{1}{K}\left( \frac{\overrightarrow{v}\cdot \overrightarrow{R}}{c} \right)-\frac{R}{K} \right\}=-1\]

\[\frac{1}{e}\left( \frac{\partial {{E}_{x}}}{\partial x}+\frac{\partial {{E}_{y}}}{\partial y}+\frac{\partial {{E}_{z}}}{\partial z} \right)=\frac{3}{{{K}^{3}}}\left( 1+\frac{\overrightarrow{a}\cdot \overrightarrow{R}}{{{c}^{2}}}-\frac{\overrightarrow{v}\cdot \overrightarrow{v}}{{{c}^{2}}} \right)\left[ \frac{R}{K}\frac{\overrightarrow{a}\cdot \overrightarrow{R}}{{{c}^{2}}}-\frac{R}{K}\frac{{{v}^{2}}}{{{c}^{2}}}-\frac{R}{K}\frac{\overrightarrow{R}\cdot \overrightarrow{a}-{{v}^{2}}}{{{c}^{2}}} \right]\]

\[\frac{1}{e}\left( \frac{\partial {{E}_{x}}}{\partial x}+\frac{\partial {{E}_{y}}}{\partial y}+\frac{\partial {{E}_{z}}}{\partial z} \right)=\frac{3}{{{K}^{3}}}\frac{R}{K}\left( 1+\frac{\overrightarrow{a}\cdot \overrightarrow{R}}{{{c}^{2}}}-\frac{\overrightarrow{v}\cdot \overrightarrow{v}}{{{c}^{2}}} \right)\left[ \frac{\overrightarrow{a}\cdot \overrightarrow{R}}{{{c}^{2}}}-\frac{{{v}^{2}}}{{{c}^{2}}}-\frac{\overrightarrow{R}\cdot \overrightarrow{a}-{{v}^{2}}}{{{c}^{2}}} \right]\]

В обозначениях программы Дробышева \(k=\frac{K}{R}\) \(K=k\cdot R\)
\(c=1\)
\[\frac{1}{e}\left( \frac{\partial {{E}_{x}}}{\partial x}+\frac{\partial {{E}_{y}}}{\partial y}+\frac{\partial {{E}_{z}}}{\partial z} \right)=\frac{1}{{{\left( k\cdot R \right)}^{3}}}\left( 1+\overrightarrow{a}\cdot \overrightarrow{R}-\overrightarrow{v}\cdot \overrightarrow{v} \right)\left[ 3+\frac{3}{k}\overrightarrow{a}\cdot \overrightarrow{R}-\frac{3}{k}{{v}^{2}}-\frac{3}{k}\left( 1+\overrightarrow{R}\cdot \overrightarrow{a}-{{v}^{2}} \right)+\frac{3}{k\cdot R}\left( \overrightarrow{v}\cdot \overrightarrow{R} \right) \right]\]
\[\frac{1}{e}\left( \frac{\partial {{E}_{x}}}{\partial x}+\frac{\partial {{E}_{y}}}{\partial y}+\frac{\partial {{E}_{z}}}{\partial z} \right)=\frac{3}{{{\left( k\cdot R \right)}^{3}}}\left( 1+\overrightarrow{a}\cdot \overrightarrow{R}-\overrightarrow{v}\cdot \overrightarrow{v} \right)\left[ 1+\frac{1}{k}\left( \overrightarrow{a}\cdot \overrightarrow{R}-{{v}^{2}} \right)-\frac{1}{k}\left( 1+\overrightarrow{R}\cdot \overrightarrow{a}-{{v}^{2}} \right)+\frac{1}{k\cdot R}\left( \overrightarrow{v}\cdot \overrightarrow{R} \right) \right]\]
\[\frac{1}{e}\left( \frac{\partial {{E}_{x}}}{\partial x}+\frac{\partial {{E}_{y}}}{\partial y}+\frac{\partial {{E}_{z}}}{\partial z} \right)=\frac{3}{{{\left( k\cdot R \right)}^{3}}}\left( 1+\overrightarrow{a}\cdot \overrightarrow{R}-\overrightarrow{v}\cdot \overrightarrow{v} \right)\left[ 1-\frac{1}{k}+\frac{1}{k\cdot R}\left( \overrightarrow{v}\cdot \overrightarrow{R} \right) \right]\]
Таким образом возникновение потенциальных ям скалярного потенциала

действительно не связано с дополнительными зарядами, материализующимися
"из воздуха", и уравнение \(\nabla\cdot\vec{E}=0\) для вакуума
выполняется. Однако для самого точечного заряда выражение дивергенции
электрического поля приходит к неопределенности вида ноль делить на
ноль. Поэтому представляет интерес произвести вычисления дивергенции
поля для неточечного заряда, имеющего заданное распределение объёмной
плотности заряда в пространстве. Интересно показать распределение
отношения радиуса Лиенара Вихерта к запаздывающему радиусу для той же
задачи

Расчёт суммарного электрического поля в тот же момент времени в целом
повторяет результат, полученный Дробышевым и Меандром в той же теме

Однако не были показаны слагаемые выражения электрического поля через
потенциалы. Первое \[\overrightarrow{{{E}_{1}}}=-\nabla \varphi \]

И второе
\[\overrightarrow{{{E}_{2}}}=-\frac{1}{c}\frac{\partial \overrightarrow{A}}{\partial t}\]

    \begin{tcolorbox}[breakable, size=fbox, boxrule=1pt, pad at break*=1mm,colback=cellbackground, colframe=cellborder]
\prompt{In}{incolor}{ }{\boxspacing}
\begin{Verbatim}[commandchars=\\\{\}]

\end{Verbatim}
\end{tcolorbox}


    % Add a bibliography block to the postdoc
    
    
    
\end{document}
